\chapter{机会约束、POMDP 与风险敏感规划}
\label{ch:plan-chance}

% ════════════════════════════════════════════════════════════════
%  不确定性下的规划与决策：把"硬保证"与"期望最优"之间的中间地带讲清。
%  融合三份 Tutorial 笔记：机会约束规划(CC-MPC/IRA/CC-RRT/DRO)、
%  POMDP 与信念空间规划(belief/α-vector/SARSOP/DESPOT/Active SLAM)、
%  风险敏感规划(一致性风险测度/CVaR/嵌套 CVaR/分布式 RL/CVaR-CBF)。
%  跨部去重：CBF/HJ 可达/鲁棒 MPC 的最深推导留待【控制部】，此处侧重规划/决策用途。
%  记号归一化：旋转 R∈SO(3)、位姿 T∈SE(3)；规划/控制自然记号保留，
%  与 planning_intro / control_intro 一致；扰动右扰动为主。
%  教学层遵循 v5.0 理论教学规范。来源 md 由 AI 生成、经多轮审，存疑处打 \pz/\rebuilt。
% ════════════════════════════════════════════════════════════════

\begin{practice}[前置自测]
开始之前，先试答下列问题。若已能流畅作答，可只速读本章表格与速查卡；若卡住，正是本章要为你解决的。
\begin{enumerate}
  \item 高斯随机变量 $X\sim\mathcal{N}(\mu,\sigma^2)$，要求 $\Pr(X\le b)\ge 1-\delta$，等价于关于 $\mu$ 的什么\emph{确定性}不等式？$\Phi^{-1}(1-\delta)$ 在其中扮演什么角色？
  \item 一条 $N$ 步轨迹，若令\emph{每一步}的碰撞概率都不超过 $\delta$，那么\emph{整条轨迹}至少碰撞一次的概率大约是多少？它等于 $\delta$ 吗？
  \item 写出马尔可夫决策过程（MDP）的五元组与 Bellman 最优方程。当\emph{状态本身不可直接观测}时，决策应当依据什么？
  \item 你用 EKF/粒子滤波估出一个状态分布。SLAM 到此为止（取均值去用），但若要在这个分布之上\emph{做决策}，还缺什么？
  \item VaR（分位点）与 CVaR（尾部平均）有何不同？两个方案碰撞概率都是 $4.9\%<5\%$，但一个轻碰、一个重创——哪种风险度量能区分它们？
\end{enumerate}
\end{practice}

\section*{本章目标}
学完本章，你应当能够：
\begin{enumerate}
  \item \textbf{形式化}机会约束（chance constraint）$\Pr(g\le 0)\ge 1-\delta$，区分单步约束 ICC 与联合约束 JCC，并用 Boole 不等式 + 风险分配把难算的 JCC 化为一族好算的 ICC；
  \item \textbf{推导}解析高斯机会约束，说明它如何等价于一个确定性的约束收紧 $\Phi^{-1}(1-\delta)\sqrt{a^\top\Sigma a}$，并落到二阶锥规划（SOCP/QP）；
  \item \textbf{掌握}迭代风险分配 IRA 的两层结构（外层分配预算、内层解凸问题），理解它为何比均分保守度低得多，以及它与安全强化学习拉格朗日法的 primal--dual 对偶；
  \item \textbf{运用}机会约束的两类规划落地——优化型的 CC-MPC 与采样型的 CC-RRT（闭环协方差传播 + 概率剔除）；并了解分布鲁棒优化 DRO 在分布本身不确定时的角色；
  \item \textbf{写出} POMDP 七元组与 belief 的贝叶斯滤波更新，论证 belief 是历史的\emph{充分统计量}，从而 POMDP 化归 belief 空间上的 MDP；
  \item \textbf{解释} $\alpha$-vector 与 PWLC 值函数（Sondik 定理）这一精确求解地基，以及"两类灾难"（维数灾难 + 历史灾难）为何令精确解爆炸；
  \item \textbf{梳理}从精确到近似的路线：QMDP$\to$点基离线（SARSOP）$\to$在线树搜索（POMCP/DESPOT），并掌握 DESPOT 的确定化稀疏 belief 树与 output-sensitive 思想；
  \item \textbf{建立} belief 空间运动规划（LQG-MP/RRBT/FIRM）作为高斯 belief 特化、以及 Active SLAM = 信息增益作奖励的 POMDP 实例；
  \item \textbf{陈述}一致性风险测度的公理、推导 CVaR 的 Rockafellar--Uryasev 变分公式并样本化成线性规划，理解静态 CVaR 的时间不一致与嵌套 CVaR 的修复，并把 CVaR 嵌进 MPC/安全滤波。
\end{enumerate}

\subsection*{本章知识导航}
本章是一条"\emph{当世界不确定、甚至看不清时，如何规划与决策}"的主线。三大支柱回答三个递进的问题：扰动是\emph{概率}的怎么办（机会约束）、状态\emph{看不清}怎么办（POMDP）、不仅要少出事还要\emph{出事别太惨}怎么办（风险敏感）。结构如下：

\begin{center}
\small
\begin{tikzpicture}[
  node distance=4mm and 6mm, >=Stealth,
  blk/.style={draw, rounded corners, align=center, font=\small, inner sep=3pt, fill=main!6, draw=main!60!black},
  grp/.style={draw, rounded corners, align=center, font=\small, inner sep=3pt, fill=second!6, draw=second!60!black},
  opt/.style={draw, rounded corners, align=center, font=\small, inner sep=3pt, fill=third!8, draw=third!60!black},
  ln/.style={->, thick, gray!60!black}]
\node[blk] (root) {不确定性下的\\规划与决策};
\node[grp, below left=8mm and 4mm of root] (cc) {机会约束\\ $\Pr(g\le0)\ge1-\delta$};
\node[grp, below=8mm of root] (pomdp) {POMDP\\ belief 空间决策};
\node[grp, below right=8mm and 4mm of root] (risk) {风险敏感\\ CVaR 尾部风险};
\node[blk, below=6mm of cc] (ccm) {ICC/JCC, IRA\\ CC-MPC, CC-RRT, DRO};
\node[blk, below=6mm of pomdp] (pm) {$\alpha$-vector, SARSOP\\ DESPOT, Active SLAM};
\node[opt, below=6mm of risk] (rm) {Artzner 公理, RU 公式\\ 嵌套 CVaR, CVaR-CBF};
\draw[ln] (root)--(cc); \draw[ln] (root)--(pomdp); \draw[ln] (root)--(risk);
\draw[ln] (cc)--(ccm); \draw[ln] (pomdp)--(pm); \draw[ln] (risk)--(rm);
\draw[ln] (ccm.east) -- (pm.west);
\draw[ln] (pm.east) -- (rm.west);
\end{tikzpicture}
\end{center}

\noindent\textbf{推荐路径}：\cref{sec:cc-form,sec:cc-gauss} 是机会约束的地基（ICC/JCC 与高斯收紧），任何读者必读；\cref{sec:cc-ira,sec:cc-rrt} 是其两条规划落地（优化型 IRA/CC-MPC、采样型 CC-RRT）；\cref{sec:pomdp-belief,sec:pomdp-alpha} 建立 POMDP 的概念地基（belief + 充分统计量 + $\alpha$-vector），\cref{sec:pomdp-solve,sec:pomdp-bsp} 给求解谱与 SLAM 接口；\cref{sec:risk-coherent,sec:risk-cvar} 是风险敏感的核心（公理 + CVaR 的 LP 化），\cref{sec:risk-nested,sec:risk-embed} 处理多步与嵌入安全工具。三支柱可独立阅读，但 \cref{sec:cc-ira}（机会约束 $\leftrightarrow$ Safe RL）、\cref{sec:risk-cvar}（CVaR 是机会约束的尾部升级）把它们焊在一起。

\subsection*{前置知识桥接}
本章建立在以下前置之上。\cref{ch:planning} 已给出运动规划三范式（搜索/采样/优化）与 RRT/RRT$^*$（\cref{sec:plan-rrt,sec:plan-rrtstar}）、动态规划与 Bellman 递推（\cref{thm:plan-bellman}）——本章的 CC-RRT 在 RRT 之上加协方差传播，POMDP 把 Bellman 搬到 belief 空间。\cref{ch:control} 给出 LQR（\cref{sec:ctrl-lqr}）、卡尔曼滤波与 LQG（\cref{sec:ctrl-lqe-dual}）、模型预测控制（\cref{sec:ctrl-mpc}）及其递归可行性（\cref{sec:ctrl-mpc-feasibility}）——本章的 CC-MPC 是给 MPC 的约束加概率收紧、闭环协方差传播要用 LQR 反馈、belief 的高斯演化就是卡尔曼那一套 Riccati 递归。当构型含姿态时，状态空间的距离/插值须用 \cref{ch:lie} 的李群语言（旋转 $\mathbf{R}\in\SO(3)$、位姿 $\mathbf{T}\in\SE(3)$）。概率论的高斯分布、正态分位数、凸优化（LP/QP/SOCP）是公共基础。

\subsection*{如果跳过本章会怎样}
\begin{itemize}
  \item 你想用 \cref{sec:ctrl-mpc} 的标称 MPC 处理风扰下的无人机，却发现\emph{风是无界的高斯}——"对有界扰动集 100\% 硬保证"的鲁棒做法根本不可行，硬约束无从谈起。机会约束给出出路：接受"被吹出边界的概率 $\le 1\%$"，问题立刻可解。
  \item 你用 EKF 估机器人在走廊的位置，但走廊两端长得一样，协方差其实是\emph{双峰}的。若无视不确定性、把均值当真值做决策，机器人会"自信地"走错方向（均值落在两峰中间的墙里）。POMDP 强迫你在\emph{整个 belief} 上决策，不确定大时先去"看清楚"。
  \item 你用机会约束做避障，要求碰撞概率 $\le5\%$。方案 A 一旦碰撞是轻碰、方案 B 一旦碰撞是重创，机会约束认为二者\emph{一样好}——它对"碰撞有多严重"完全失明。风险敏感（CVaR）看尾部的平均代价，明确区分 A 与 B。
\end{itemize}

\begin{table}[H]\centering\small
\caption{本章阅读方式建议}
\label{tab:plan-chance-reading}
\begin{tabular}{lll}
\toprule
阅读方式 & 时间 & 适合谁 \\
\midrule
精读（含全部推导与练习） & 8--10 小时 & 首次系统学习不确定性规划 \\
速读（跳过 DESPOT 家族与嵌套 CVaR 推导） & 3 小时 & 有 MPC/滤波背景、想对齐脉络 \\
速查（只看小结的符号表 + 定理速查表） & 20 分钟 & 写代码、选方法时回来查 \\
\bottomrule
\end{tabular}
\end{table}

\begin{note}
\textbf{本章记号与约定.}\;
风险预算（允许的违反概率）记 $\delta\in(0,1)$；置信水平 $1-\delta$。标准正态 CDF 记 $\Phi$，其逆（分位函数）$\Phi^{-1}$。状态/约束量 $\mathbf{x}$ 的均值记 $\bar{\mathbf{x}}$、协方差 $\boldsymbol{\Sigma}$ 或沿轨迹的 $\mathbf{P}_k$。POMDP 的 belief（状态后验）记 $b$，$b(s)$ 是处于状态 $s$ 的概率；为避免与机会约束的违反概率冲突，POMDP 的折扣因子记 $\gamma$、观测空间记 $\Omega$（\emph{本章 $\Omega$ 指观测空间，与控制部 RPI 不变集 $\Omega$ 无关}）。风险度量记 $\rho$，CVaR 的尾部水平记 $\alpha\in(0,1)$（$\mathrm{CVaR}_\alpha$ 看最坏 $1-\alpha$ 尾部）。\textbf{代价方向约定}：本章统一把随机变量 $Z$/$X$/$L$ 当作\emph{代价/损失、越大越坏}（与金融文献"收益越大越好"相反，引用时须翻译符号方向）。旋转 $\mathbf{R}\in\SO(3)$、位姿 $\mathbf{T}\in\SE(3)$；状态不确定性沿用右扰动直觉。$\boldsymbol{\alpha}$-vector 的字母 $\alpha$ 与 CVaR 的 $\alpha$ 同形但语境不同，已在各处注明。
\end{note}

% ════════════════════════════════════════════════════════════════
\section{机会约束的形式化：ICC、JCC 与 Boole 不等式}
\label{sec:cc-form}

\paragraph{动机：无界扰动下，"零违反"是个伪命题。}
\cref{sec:ctrl-mpc} 的标称 MPC 假设动力学精确；鲁棒做法则对一个\emph{有界}扰动集 $w\in\mathcal{W}$ 给 100\% 保证。但现实里很多扰动是\emph{无界}的——高斯风扰、传感器高斯噪声、未建模残差。对无界扰动，"零违反概率"在数学上不可能（任意大的扰动都有非零概率）\cite{blackmore2011chance}。于是问题必须换个问法：不再要求"绝不违反约束 $g(\mathbf{x})\le 0$"，而是要求"\emph{违反的概率不超过一个我能接受的小值 $\delta$}"：
\begin{equation}
  \Pr\big(g(\mathbf{x})\le 0\big)\ \ge\ 1-\delta.
  \label{eq:cc-def}
\end{equation}
这就是\textbf{机会约束}（chance constraint）。$\delta$ 是你愿意为"少保守、多性能"付出的风险代价——$\delta=0.01$ 表示允许 $1\%$ 的违反概率。

\paragraph{反面：手调安全裕量，要么太保守要么不够。}
不引入概率框架的朴素做法是给约束手动加一个固定裕量（margin）防扰动。但窄通道里固定裕量太大就过不去、放小了又偶尔违反。机会约束把裕量从"拍脑袋的常数"变成"由协方差与风险 $\delta$ 算出来的量"（\cref{sec:cc-gauss}）——该大则大、该小则小，且违反概率有形式化保证。

\begin{insight}
机会约束的精髓是把"安全"从一个\textbf{布尔判断}（撞/不撞）变成一个\textbf{连续资源}（违反概率 $\delta$）。一旦安全成了可分配的预算，你就能问出鲁棒做法问不出的问题：把有限的风险预算花在哪一步、哪个约束上最划算？这正是迭代风险分配（\cref{sec:cc-ira}）的全部动机，也是机会约束相对硬约束的根本优势来源——\textbf{用"接受一点点风险"换"少很多保守度"}。
\end{insight}

\subsection{两种形式：单步约束 ICC 与联合约束 JCC}
\label{sec:cc-icc-jcc}

\begin{definition}[单步机会约束 ICC]\label{def:cc-icc}
\textbf{单步机会约束}（Individual Chance Constraint, ICC）要求每个时步 $k$、每个约束面 $j$ \emph{独立}满足，记 $a_j^\top\mathbf{x}_k\le b_j$ 为第 $j$ 个半空间约束（安全侧），则
\begin{equation}
  \Pr\big(a_j^\top\mathbf{x}_k\le b_j\big)\ \ge\ 1-\delta_j,\qquad \forall k,\ \forall j.
  \label{eq:cc-icc}
\end{equation}
直觉（赛车例子）：ICC 限制的是"在\emph{某一个}时刻撞墙的概率"。
\end{definition}

\begin{definition}[联合机会约束 JCC]\label{def:cc-jcc}
\textbf{联合机会约束}（Joint Chance Constraint, JCC）要求\emph{整条轨迹的所有约束一起}满足：
\begin{equation}
  \Pr\Big(\bigcap_{k=0}^{N}\bigcap_{j=1}^{M} a_j^\top\mathbf{x}_k\le b_j\Big)\ \ge\ 1-\delta.
  \label{eq:cc-jcc}
\end{equation}
直觉：JCC 限制的是"\emph{从起点到终点全程}至少违反一次的概率"。
\end{definition}

\paragraph{为什么 JCC 更关键也更难。}
你真正关心的几乎总是 JCC——"这趟飞行安全到达的概率 $\ge99\%$"，而不是"每一个时刻单独看安全的概率 $\ge99\%$"。但 JCC \cref{eq:cc-jcc} 里那个对所有 $k,j$ 取交集的联合概率，一般\emph{没有解析表达}（涉及多维相关事件的联合分布），直接算不可行。ICC 则因为每个约束独立、在高斯下有解析形式（\cref{sec:cc-gauss}），好算得多。\textbf{核心技术难题就是：如何用一族好算的 ICC 来保证一个难算的 JCC。}

\subsection{Boole 不等式：把 JCC 拆成一族 ICC}
\label{sec:cc-boole}

桥梁是\textbf{Boole 不等式}（union bound）。设事件 $A_{j,k}$ 表示"第 $k$ 步第 $j$ 个约束被违反"，则任意事件的并满足
\begin{equation}
  \Pr\Big(\bigcup_{j,k} A_{j,k}\Big)\ \le\ \sum_{j,k}\Pr(A_{j,k}).
  \label{eq:cc-boole}
\end{equation}

\begin{theorem}[Boole 拆解：一族 ICC 保证 JCC]\label{thm:cc-boole}
若对每个 $(j,k)$ 令违反概率 $\Pr(A_{j,k})\le\delta_{j,k}$，且这些小风险之和满足 $\sum_{j,k}\delta_{j,k}\le\delta$，则 JCC \cref{eq:cc-jcc} 成立。
\end{theorem}
\begin{proof}
"全程至少违反一次"的概率即 $\Pr(\bigcup_{j,k}A_{j,k})$。由 \cref{eq:cc-boole} 及假设，
\[
  \Pr(\text{全程违反})\ \le\ \sum_{j,k}\Pr(A_{j,k})\ \le\ \sum_{j,k}\delta_{j,k}\ \le\ \delta,
\]
故 $\Pr(\text{全程不违反})\ge 1-\delta$。
\end{proof}

\cref{thm:cc-boole} 把"难算的联合概率"换成了"好算的一族单步概率 + 一个求和约束"——只要每个 ICC 用各自的小风险 $\delta_{j,k}$、且这些小风险加起来不超过总预算 $\delta$，JCC 自动成立。

\begin{note}
\textbf{Boole 不等式是保守的.}\;union bound 只在各违反事件\emph{互斥}时取等；实际中它们常重叠（同一段轨迹可能在相邻时刻同时接近同一面墙），所以 $\sum\delta_{j,k}$ 高估了真实联合违反概率——这意味着用 Boole 拆出来的 JCC 解\emph{偏保守}（真实违反概率往往远小于 $\delta$）。这点保守度是"把难算问题变好算"的代价。
\end{note}

\paragraph{风险分配：$\delta_{j,k}$ 怎么分？}
\cref{thm:cc-boole} 只要求 $\sum_{j,k}\delta_{j,k}\le\delta$，但\emph{怎么把总预算 $\delta$ 分给各个 $\delta_{j,k}$}，大有讲究：
\begin{itemize}
  \item \textbf{均分}（uniform allocation）：$\delta_{j,k}=\delta/(NM)$。最简单，但很保守——它给"离障碍很远、几乎不可能违反"的约束和"贴着障碍、很危险"的约束分了\emph{一样多}的预算，浪费。
  \item \textbf{最优分配}（IRA，\cref{sec:cc-ira}）：把预算从"过度宽松"的约束（距离 $\gg$ margin）转移给"紧绷"的约束。直觉：风险预算是稀缺资源，该花在刀刃（紧约束）上。
\end{itemize}

\begin{pitfall}[概念误区：把 ICC 当成 JCC 用]
新手常以为"我对每一步都加了 $\Pr(\text{违反})\le\delta$，那整条轨迹的违反概率也是 $\delta$ 吧？"。实际整条轨迹的违反概率远大于 $\delta$——若每步独立违反概率都是 $\delta$、共 $N$ 步，全程至少违反一次的概率接近 $1-(1-\delta)^N\approx N\delta$，是单步的 $N$ 倍。根因：ICC 管单步、JCC 管全程，两者差一个 union bound。要 JCC $=\delta$，每个 ICC 应设为 $\delta_{j,k}$ 且 $\sum\delta_{j,k}\le\delta$。检验方法：蒙特卡洛统计"整条轨迹"的违反率，而非"单步"违反率。
\end{pitfall}

\begin{practice}
\begin{enumerate}
  \item 对 1D 高斯 $X\sim\mathcal{N}(0,1)$，分别算 $\delta=0.1,0.05,0.01$ 对应的 $\Phi^{-1}(1-\delta)$（即约束收紧系数）。验证 $\delta$ 越小、系数越大、收紧越多。
  \item 一条 20 步轨迹、每步 1 个约束，要求 JCC 总违反概率 $\le0.05$。若均分，每步 ICC 的 $\delta_k$ 是多少？对应的 $\Phi^{-1}(1-\delta_k)$ 比 JCC 直接用 $0.05$ 大多少？这个差距就是 Boole + 均分的保守度代价。
  \item 构造一个 2D 高斯随机游走轨迹，分别按"每步 ICC $=0.05$"和"整条 JCC $=0.05$（均分）"设约束，蒙特卡洛对比两者的\emph{整条轨迹}违反率，直观感受 ICC $\neq$ JCC。
\end{enumerate}
\end{practice}

% ════════════════════════════════════════════════════════════════
\section{解析高斯收紧与机会约束 MPC}
\label{sec:cc-gauss}

\paragraph{动机：一般机会约束 NP-hard，必须找可解的近似。}
坏消息先说：\emph{一般的机会约束优化是 NP-hard 的}\cite{nemirovski2006convex,calafiore2006distributionally}——可行域 $\{\mathbf{x}:\Pr(g(\mathbf{x})\le0)\ge1-\delta\}$ 一般\emph{非凸}（即使每个 $g$ 凸、扰动良性，"以概率 $1-\delta$ 满足"这个要求也可能切出非凸甚至不连通的可行集，因为概率约束本质是对"违反集合的测度"设限，而测度对决策变量的依赖可高度非线性）。好消息是：在特定假设（如高斯 + 线性约束）下，机会约束可化成可高效求解的凸问题。

\subsection{解析高斯机会约束的完整推导}
\label{sec:cc-gauss-derive}

\begin{theorem}[解析高斯机会约束的确定性等价]\label{thm:cc-gauss}
设状态（或约束量）$\mathbf{x}\sim\mathcal{N}(\bar{\mathbf{x}},\boldsymbol{\Sigma})$，线性约束 $a^\top\mathbf{x}\le b$。则机会约束 $\Pr(a^\top\mathbf{x}\le b)\ge1-\delta$ 等价于确定性约束
\begin{equation}
  a^\top\bar{\mathbf{x}}\ \le\ b-\Phi^{-1}(1-\delta)\,\sqrt{a^\top\boldsymbol{\Sigma}a}.
  \label{eq:cc-gauss-tighten}
\end{equation}
\end{theorem}

\begin{derivation}
分四步\cite{blackmore2011chance}。
\emph{(1) 线性变换保持高斯}：$y\triangleq a^\top\mathbf{x}$ 是高斯，$y\sim\mathcal{N}(a^\top\bar{\mathbf{x}},\,a^\top\boldsymbol{\Sigma}a)$。记标准差 $s=\sqrt{a^\top\boldsymbol{\Sigma}a}$。
\emph{(2) 标准化}：
\[
  \Pr(y\le b)=\Pr\!\Big(\tfrac{y-a^\top\bar{\mathbf{x}}}{s}\le\tfrac{b-a^\top\bar{\mathbf{x}}}{s}\Big)=\Phi\!\Big(\tfrac{b-a^\top\bar{\mathbf{x}}}{s}\Big),
\]
其中 $\Phi$ 是标准正态 CDF。
\emph{(3) 施加约束}：要 $\Phi\big(\tfrac{b-a^\top\bar{\mathbf{x}}}{s}\big)\ge1-\delta$，两边取单调的 $\Phi^{-1}$：$\tfrac{b-a^\top\bar{\mathbf{x}}}{s}\ge\Phi^{-1}(1-\delta)$。
\emph{(4) 整理}：移项即得 \cref{eq:cc-gauss-tighten}。
\end{derivation}

\paragraph{解读。}
原约束的均值版 $a^\top\bar{\mathbf{x}}\le b$ 被\emph{向内收紧}了 $\Phi^{-1}(1-\delta)\sqrt{a^\top\boldsymbol{\Sigma}a}$。收紧量由两部分构成：风险因子 $\Phi^{-1}(1-\delta)$（$\delta$ 越小越大；$\delta=0.05\Rightarrow1.645$，$\delta=0.01\Rightarrow2.326$）乘以约束方向上的标准差 $\sqrt{a^\top\boldsymbol{\Sigma}a}$（不确定性越大越大）。

\paragraph{为什么这是 SOCP。}
$\sqrt{a^\top\boldsymbol{\Sigma}a}=\|\boldsymbol{\Sigma}^{1/2}a\|_2$ 是一个二范数。把决策变量（如均值轨迹）代入，约束形如 $a^\top\bar{\mathbf{x}}+\Phi^{-1}(1-\delta)\|\boldsymbol{\Sigma}^{1/2}a\|_2\le b$——左边是"线性项 + 二范数项"，正是\textbf{二阶锥（SOC）约束}的标准形。所以解析高斯 CC-MPC 是一个 SOCP（若 $\boldsymbol{\Sigma}$ 固定、只优化均值且约束线性，常可进一步简化为带锥约束的 QP）。

\begin{insight}
整个解析高斯机会约束的魔法只有一句话：\textbf{概率约束 $\Pr(\cdot)\ge1-\delta$ 等价于在确定性约束上加一个正比于"标准差 $\times$ 分位数"的安全 margin}。$\delta$ 控制分位数（要多稳），$\boldsymbol{\Sigma}$ 控制标准差（有多不确定），两者相乘就是该收多少。三个视角看同一公式：\emph{几何上}——把约束面（墙）沿法向向内平移 $\Phi^{-1}(1-\delta)\sqrt{a^\top\boldsymbol{\Sigma}a}$；\emph{概率上}——margin 恰好罩住高斯在该方向上 $1-\delta$ 的概率质量；\emph{优化上}——它是一个凸的二阶锥约束。CC-MPC、CC-RRT、协方差控制的约束处理，都在反复用这一个收紧。
\end{insight}

\subsection{机会约束 MPC 与闭环协方差传播}
\label{sec:cc-mpc}

把 \cref{thm:cc-gauss} 用到整条 MPC 轨迹，即得\textbf{机会约束 MPC}（CC-MPC）。考虑线性系统 $\mathbf{x}_{k+1}=\mathbf{A}\mathbf{x}_k+\mathbf{B}\mathbf{u}_k+\mathbf{w}_k$，过程噪声 $\mathbf{w}_k\sim\mathcal{N}(\mathbf{0},\boldsymbol{\Sigma}_w)$，障碍用半空间 $a_j^\top\mathbf{x}\le b_j$ 描述。关键的一步是\textbf{闭环协方差传播}：在反馈控制 $\mathbf{u}_k=\bar{\mathbf{u}}_k-\mathbf{K}(\mathbf{x}_k-\bar{\mathbf{x}}_k)$ 下，误差服从闭环动力学，状态协方差按
\begin{equation}
  \mathbf{P}_{k+1}=(\mathbf{A}-\mathbf{B}\mathbf{K})\,\mathbf{P}_k\,(\mathbf{A}-\mathbf{B}\mathbf{K})^\top+\boldsymbol{\Sigma}_w
  \label{eq:cc-cov-prop}
\end{equation}
逐步演化（这是离散李雅普诺夫递推；增益 $\mathbf{K}$ 可取 \cref{sec:ctrl-lqr} 的 LQR 增益 \cref{eq:ctrl-dlqr-K}）。注意 $\mathbf{P}_k$ \emph{与控制序列无关}，可\emph{离线预先算好}——这是解析高斯 CC-MPC 高效的关键：在线只剩一个 QP。每个障碍面用 $\Phi^{-1}(1-\delta_{j,k})\sqrt{a_j^\top\mathbf{P}_k a_j}$ 向内收紧，再用 Boole 均分（或 IRA）把 JCC 拆到每步每面。

\begin{note}
\textbf{必须用闭环、不能用开环传播.}\;若用开环 $\mathbf{P}_{k+1}=\mathbf{A}\mathbf{P}_k\mathbf{A}^\top+\boldsymbol{\Sigma}_w$，在 $\mathbf{A}$ 不稳定时协方差\emph{发散}，收紧 margin 巨大、问题不可行。反馈把偏差/协方差镇定住——这与 \cref{sec:ctrl-mpc} 的 tube 思想一脉相承（用反馈构造不变集）。\pz{控制部深讲：tube MPC 的鲁棒正不变集 (RPI) 收紧与本节闭环协方差收紧的精确联系，待 \cref{sec:ctrl-mpc-feasibility} 后续鲁棒 MPC 章。}
\end{note}

\begin{codebox}[Python]{解析 CC-MPC 的约束收紧（核心）}
import numpy as np, cvxpy as cp
from scipy.stats import norm

def cc_mpc(x0, goal, obstacles, A, B, Sigma_w, K, N=20,
           delta=0.05, u_max=3.0):
    # obstacles: [(a_j, b_j)] safe side a_j^T p <= b_j ; p = x[:2]
    nx, nu = B.shape
    M = max(len(obstacles), 1)
    delta_step = delta / (N * M)            # Boole uniform allocation
    z = norm.ppf(1.0 - delta_step)          # Phi^{-1}(1 - delta_step)
    A_cl = A - B @ K                        # closed-loop matrix
    # 1) propagate closed-loop covariance offline (control-independent)
    P = [np.zeros((nx, nx))]
    for _ in range(N):
        P.append(A_cl @ P[-1] @ A_cl.T + Sigma_w)
    # 2) online QP over nominal (mean) trajectory z_var, control v_var
    z_var = cp.Variable((nx, N + 1)); v_var = cp.Variable((nu, N))
    xg = np.concatenate([goal, np.zeros(nx - 2)])
    Q = np.diag([5., 5., 1., 1.]); R = 0.1 * np.eye(nu)
    cost, cons = 0, [z_var[:, 0] == x0]
    for k in range(N):
        cost += cp.quad_form(z_var[:, k] - xg, Q) + cp.quad_form(v_var[:, k], R)
        cons += [z_var[:, k+1] == A @ z_var[:, k] + B @ v_var[:, k]]
        cons += [cp.abs(v_var[:, k]) <= u_max]
        Ppos = P[k][:2, :2]                 # position-block covariance
        for a_j, b_j in obstacles:
            margin = z * np.sqrt(float(a_j @ Ppos @ a_j))
            cons += [a_j @ z_var[:2, k] <= b_j - margin]   # inward tightening
    cp.Problem(cp.Minimize(cost), cons).solve(solver=cp.OSQP, warm_start=True)
    return v_var[:, 0].value
\end{codebox}

\paragraph{数值走查（可手算验证）。}
用一维例子把收紧算清。单约束 $p\le1$（离墙 1 米），位置标准差 $\sigma=0.2$（来自传播的协方差），风险 $\delta=0.05$：收紧系数 $\Phi^{-1}(0.95)\approx1.645$，收紧量 $1.645\times0.2=0.329$，收紧后名义需满足 $\bar p\le1-0.329=0.671$。即为把"撞墙概率 $\le5\%$"，名义轨迹必须离墙 $0.329$ 米以上。若 $\delta=0.01$（$\Phi^{-1}(0.99)\approx2.326$），收紧量增至 $0.465$、$\bar p\le0.535$，更保守；若 $\sigma$ 减半，收紧量也减半至 $0.165$——\textbf{margin 随不确定性线性缩放}。\pz{存疑：源称 $\delta=0.01\Rightarrow\Phi^{-1}=2.326$、$\delta=0.05\Rightarrow1.645$，这两个分位数值标准、可信；但"$\Phi^{-1}(0.95)\approx1.642$/经验 $1.645$"等小数末位在源不同处略有出入，取教科书值 $1.645$、$2.326$。}

\subsection{机会约束近似方法谱}
\label{sec:cc-spectrum}

解析高斯只是一种近似。\cref{tab:cc-spectrum} 按"分布假设强弱"排列完整方法谱。

\begin{table}[H]\centering\small
\caption{机会约束近似方法谱（按分布假设强度排列）}
\label{tab:cc-spectrum}
\begin{tabular}{lllll}
\toprule
方法 & 分布假设 & 求解类型 & 保守性 & 代表工作 \\
\midrule
解析高斯 $\Phi^{-1}$ & 高斯 & SOCP/QP & ICC 精确 & Blackmore--Ono--Williams\cite{blackmore2011chance} \\
Cantelli / 椭球 & 仅知 $\mu,\boldsymbol{\Sigma}$ & SOCP & 保守 & Calafiore--El Ghaoui\cite{calafiore2006distributionally} \\
Boole + IRA & 高斯/矩 & 外层 NLP + 内层 SOCP & 可收敛 & Ono--Williams\cite{ono2008iterative} \\
粒子 / 场景 MILP & 任意 & MILP & 样本越多越准 & Blackmore et al.\cite{blackmore2010probabilistic} \\
场景法（scenario） & 分布无关 & LP/QP & 由样本数控制 & Campi--Garatti\cite{campi2008exact} \\
DR-CC（Wasserstein/矩）& 分布族 & SDP/SOCP & 最保守（数据少最稳）& \cite{calafiore2006distributionally} \\
\bottomrule
\end{tabular}
\end{table}

逐个要点：
\begin{itemize}
  \item \textbf{Cantelli（单边切比雪夫）}：只知均值与协方差、不知分布形状时，用收紧系数 $\sqrt{(1-\delta)/\delta}$ 代替 $\Phi^{-1}(1-\delta)$。无需高斯假设，但更保守（如 $\delta=0.05$：Cantelli 系数 $\approx4.36$，高斯仅 $1.645$）。\pz{存疑：Cantelli 单边界为 $\Pr(X-\mu\ge t\sigma)\le 1/(1+t^2)$，令右端 $=\delta$ 得 $t=\sqrt{(1-\delta)/\delta}$；$\delta=0.05$ 时 $t=\sqrt{0.95/0.05}=\sqrt{19}\approx4.36$，与源一致，公式与数值核对无误。}
  \item \textbf{粒子 / 场景 MILP}：用大量采样的扰动粒子，把"违反粒子比例 $\le\delta$"写成 MILP（二元变量标记每个粒子是否违反）。可处理任意分布，但 MILP 随粒子数/时步爆炸。
  \item \textbf{场景法}（Campi--Garatti\cite{campi2008exact}）：用 $N$ 个独立采样的约束替换概率约束，给出\emph{分布无关}的样本复杂度界——大致 $N\ge\frac{2}{\delta}(d+\ln\frac1\beta)$（$d$ 为决策维、$\beta$ 为置信），采样数足够大时解以高概率满足原机会约束。优雅且通用，代价是 $N$ 可能很大。这里的 $N\ge\lceil\frac{2}{\delta}(d+\ln\frac1\beta)\rceil$ 是经典的 Calafiore--Campi / Campi--Garatti--Prandini 界（取整后），已核对无误；注意它对应"约束不可丢弃"的基本版本，sampling-and-discarding / wait-and-judge 等变体的常数与形式不同。
  \item \textbf{分布鲁棒机会约束（DR-CC）}：连分布都不确定时，对"一个分布族（矩约束 / Wasserstein 球）内的最坏分布"保证机会约束。最稳健（尤其数据少时），化成 SDP/SOCP。这是分布鲁棒优化（DRO）思想在机会约束上的落地。
\end{itemize}

\paragraph{工程选择。}
高频 MPC 首选"解析高斯 + IRA"（快 + 准）；分布明显非高斯 $\to$ 粒子/场景 MILP 或场景法；有限数据 + 分布模糊 $\to$ Wasserstein DR-CC。

\begin{pitfall}[概念误区：Cantelli 与高斯收紧混用]
"都是收紧 margin，用哪个系数无所谓"——错。$\Phi^{-1}(1-\delta)$ 依赖高斯假设；Cantelli 系数 $\sqrt{(1-\delta)/\delta}$ 只用均值方差、对任意分布成立但保守。明明是高斯却用 Cantelli（系数 $4.36$ vs $1.645$），过度保守；明明非高斯却用高斯系数，违反概率超标。正确做法：确认分布——高斯（或近似高斯）用 $\Phi^{-1}$；只知矩、分布未知用 Cantelli；明显非高斯用粒子/场景法。
\end{pitfall}

\begin{practice}
\begin{enumerate}
  \item 对 2D 双积分器 + 高斯过程噪声实现解析 CC-MPC，每个障碍面用 $\Phi^{-1}(1-\delta)$ margin 收紧，OSQP 求解。蒙特卡洛 1000 次验证碰撞率 $\le\delta$；对比 $\delta=0.01$ 与 $\delta=0.1$ 的轨迹保守度差异。
  \item 对同一机会约束问题分别用解析高斯、Cantelli、场景法（采样）三种近似，对比保守度（实际违反率距 $\delta$ 多远）与求解时间。
  \item 论证：为什么一般问题非凸、但高斯 + 线性约束下变凸？（提示：$\Phi^{-1}(1-\delta)\sqrt{a^\top\boldsymbol{\Sigma}a}$ 关于均值线性、关于决策变量是二阶锥的，凸性来自这里；一般分布/非线性约束破坏这个结构。）
\end{enumerate}
\end{practice}

% ════════════════════════════════════════════════════════════════
\section{迭代风险分配 IRA 与 Safe RL 的对偶}
\label{sec:cc-ira}

\paragraph{动机：均分浪费预算。}
\cref{sec:cc-boole} 说均分风险很保守——它给离障碍很远的约束和贴着障碍的约束分了一样多的预算。\textbf{迭代风险分配}（Iterative Risk Allocation, IRA；Ono--Williams\cite{ono2008iterative}）用一个两阶段迭代把风险预算分到刀刃上。

\subsection{IRA 的两层结构}
\label{sec:cc-ira-struct}

\begin{itemize}
  \item \textbf{下层}（lower stage）：给定当前风险分配 $\{\delta_{j,k}\}$，每个 ICC 收紧成确定性约束（解析高斯 margin），解一个\emph{凸问题}（LP/QP/SOCP）得最优均值轨迹。
  \item \textbf{上层}（upper stage）：固定下层轨迹，\emph{重新分配}风险——把"不活跃约束"（实际距离 $\gg$ 收紧 margin，即风险没用满）的多余预算，转移给"活跃约束"（贴着边界、margin 吃紧）。上层目标凸（但不总可微），用快速下降算法求解。
  \item \textbf{迭代}：下层、上层交替，直到风险分配收敛、总代价不再下降。
\end{itemize}

\paragraph{为什么有效。}
均分把预算平摊给所有约束，但大部分约束根本用不到那么多（离障碍远）。IRA 把这些浪费的预算挪给紧约束，让紧约束能用更大的 $\delta_{j,k}$（更小的收紧 margin），于是整体可行域更大、轨迹更优。Ono--Williams\cite{ono2008iterative} 证明：IRA 的次优性远小于椭球松弛法，同时比粒子法（MILP）快得多——它卡在"保守度"与"计算量"权衡的一个甜点上。

\begin{insight}
IRA $=$ 风险预算的"按需分配"。把违反概率 $\delta$ 想成一笔钱：均分是"人人平均发"，IRA 是"按需分配——谁紧给谁多"。数学上，IRA 是在 ICC 可行性约束下最小化总代价、对风险分配 $\{\delta_{j,k}\}$ 做优化。
\end{insight}

\begin{codebox}[Python]{IRA 外层骨架（教学版）}
import numpy as np

def ira(solve_lower, active_margin, delta_total, NM, n_iter=20, tol=1e-4):
    # solve_lower(deltas) -> (cost, traj, slack):
    #   slack = actual_distance - tightened_margin (>0 means inactive)
    deltas = np.full(NM, delta_total / NM)         # init: uniform
    prev_cost = np.inf
    for _ in range(n_iter):
        cost, traj, slack = solve_lower(deltas)    # lower: convex problem
        active = slack <= active_margin            # tight constraints
        n_active = max(int(active.sum()), 1)
        # upper: move surplus risk from inactive to active constraints
        residual = deltas[~active].sum() * 0.5     # reclaim half (damped)
        deltas[~active] *= 0.5
        deltas[active]  += residual / n_active
        deltas = np.clip(deltas, 1e-6, None)
        deltas *= delta_total / deltas.sum()       # renormalize: sum(delta) = delta_total
        if abs(prev_cost - cost) < tol:
            break
        prev_cost = cost
    return deltas, traj
\end{codebox}

\begin{pitfall}[编程陷阱：忘记重新归一化导致风险预算超支]
IRA 迭代几轮后，$\sum\delta_{j,k}$ 悄悄超过 $\delta$（风险转移时若只加不减、或数值误差累积），JCC 保证失效，蒙特卡洛违反率超标。正确写法：每轮迭代后强制 \texttt{deltas *= delta\_total / deltas.sum()} 重新归一化，确保 $\sum\delta_{j,k}=\delta$ 始终成立。检验：每轮打印 \texttt{deltas.sum()}，应恒等于 $\delta$。
\end{pitfall}

\begin{pitfall}[思维陷阱：以为 IRA 能消除 Boole 的保守度]
"IRA 优化了风险分配，那 Boole 的保守度就没了"——错。IRA 优化的是"在 Boole 框架内\emph{怎么分}预算"，但 Boole 的 union bound 上界本身的保守度（事件重叠导致）IRA 消不掉。IRA 让分配最优，不让 Boole 变精确——两者保守度叠加。要进一步减后者需更紧的联合风险界。
\end{pitfall}

\subsection{机会约束与安全强化学习的 primal--dual 对偶}
\label{sec:cc-saferl}

做强化学习的读者会发现：CC-MPC 与\textbf{安全强化学习}（约束马尔可夫决策过程，CMDP）长得很像——都在"最大化性能"的同时"限制违反"。本节点明：\emph{IRA 的风险分配与 Safe RL 的拉格朗日乘子调节，是同一个优化问题的 primal--dual 两端}。

\begin{table}[H]\centering\small
\caption{CC-MPC 与 Safe RL（CMDP）的对照}
\label{tab:cc-saferl}
\begin{tabular}{lll}
\toprule
维度 & CC-MPC & Safe RL（CMDP）\\
\midrule
模型 & 已知动力学 $f$、噪声 $\boldsymbol{\Sigma}$ & 未知 MDP，靠样本学习 \\
约束 & $\Pr(g\le0)\ge1-\delta$（机会约束）& $\mathbb{E}[\sum_t\gamma^t c_t]\le d$（期望累积代价约束）\\
求解 & SDP / SOCP / MILP / SCP & Trust-region / Primal-dual \\
代表 & Ono--Williams IRA\cite{ono2008iterative} & CPO\cite{achiam2017constrained}、PPO-Lagrangian \\
\bottomrule
\end{tabular}
\end{table}

\paragraph{统一拉格朗日。}
两者都可写成"性能目标 + 约束惩罚"的拉格朗日：
\begin{equation}
  \mathcal{L}(\pi,\lambda)=J(\pi)+\sum_i\lambda_i\big(C_i(\pi)-d_i\big),
  \label{eq:cc-lagrangian}
\end{equation}
其中 $J(\pi)$ 是性能/代价，$C_i$ 是第 $i$ 个约束的违反度量、$d_i$ 是其预算、$\lambda_i\ge0$ 是对应乘子。优化在 $\min_\pi\max_{\lambda\ge0}\mathcal{L}$ 的鞍点。区别只在\emph{从哪一端逼近这个鞍点}：
\begin{itemize}
  \item \textbf{CC-MPC / IRA：从 primal 端（分配预算 $d_i$）}。IRA 直接调整各约束的风险预算 $\delta_{j,k}$（即 $d_i$），下层在给定预算下解凸问题。它在\emph{原始空间}操作"每个约束分多少违反额度"。
  \item \textbf{Safe RL / PPO-Lagrangian：从 dual 端（调乘子 $\lambda_i$）}。不直接分配预算，而是自适应调节乘子——约束违反了就增大 $\lambda_i$（加重惩罚）、满足了就减小。它在\emph{对偶空间}操作"每个约束的惩罚价格"。
\end{itemize}

\begin{insight}
风险是预算，也是价格。CC 把违反概率当\emph{预算}（primal，分配 $\delta$），Safe RL 把违反当\emph{价格}（dual，调乘子 $\lambda$）。风险预算 $\delta_i$（primal）与乘子 $\lambda_i$（dual）是一对对偶变量——分配更多预算 $\Longleftrightarrow$ 更低的影子价格（乘子）。IRA"把预算挪给紧约束"与 Lagrangian"给紧约束更高的 $\lambda$"，在鞍点处描述的是\emph{同一个最优}，只是一个从量、一个从价逼近。掌握这个对偶，你就能在"已知模型的优化"与"未知模型的学习"之间自由翻译。
\end{insight}

\begin{derivation}
用一个可手算的标量例子把对偶说实。把"风险预算"抽象成一般资源预算 $d$，问题为 $\min_x(x-2)^2$ s.t. $x\le d$。
\emph{Primal（给定 $d$）}：若 $d\ge2$，约束不起作用，$x^\star=2$、$\lambda^\star=0$；若 $d<2$，约束起作用，$x^\star=d$、$J^\star=(d-2)^2$，预算的影子价格 $\lambda^\star=-\frac{\mathrm{d}J^\star}{\mathrm{d}d}=2(2-d)>0$（预算越紧、价格越高）。
\emph{Dual（调 $\lambda$）}：拉格朗日 $\mathcal{L}(x,\lambda)=(x-2)^2+\lambda(x-d)$，对 $x$ 求极小得 $x=2-\lambda/2$；令约束活跃（$d<2$ 时 $x=d$）：$2-\lambda/2=d\Rightarrow\lambda=2(2-d)$——\emph{与 primal 的影子价格完全相同}。映射到机会约束：把 $d$ 换成风险预算 $\delta_i$、约束换成"违反概率 $\le\delta_i$"，IRA 调 $\delta_i$（primal）、PPO-Lagrangian 调 $\lambda_i$（dual），两者在鞍点描述同一最优。
\end{derivation}

\begin{practice}
\begin{enumerate}
  \item 在 \cref{sec:cc-mpc} 的 CC-MPC 上实现 IRA（外层分配 $\delta_{j,k}$、内层解 QP/SOCP）。对比 IRA 与均分在同一窄通道任务上的代价与可行性，定量说明 IRA 降了多少保守度。
  \item 用一个简单的约束 LQR/CMDP 例子，分别用"固定预算解凸问题（primal）"和"自适应乘子梯度上升（dual）"求解，验证两者收敛到同一最优，观察 $\delta$ 与 $\lambda$ 的对偶关系。
  \item 写出 CC-MPC 与 Safe RL 统一的拉格朗日 \cref{eq:cc-lagrangian}，说明 IRA 与 PPO-Lagrangian 各在哪个空间、调什么变量，并解释为什么风险预算 $\delta$ 与乘子 $\lambda$ 是对偶变量。
\end{enumerate}
\end{practice}

% ════════════════════════════════════════════════════════════════
\section{机会约束的采样型落地：CC-RRT 与协方差控制}
\label{sec:cc-rrt}

\paragraph{动机：让采样规划器也懂不确定性。}
前面的机会约束都绑在\emph{优化型}规划器（MPC）里。但 \cref{ch:planning} 介绍的另一大类是\emph{采样型}规划器——RRT、RRT$^*$、PRM。它们假设状态精确、障碍精确，真机一有扰动/定位误差，规划出的路径就可能不可行。\textbf{CC-RRT}（Luders--Kothari--How\cite{luders2010chance}）把机会约束嵌进 RRT：树生长时沿每条边\emph{传播状态协方差}，按"违反概率 $\le\delta$"\emph{剔除危险节点}，从而让采样规划器输出的路径带"以 $1-\delta$ 概率安全"的保证。

\subsection{CC-RRT 的三个核心机制}
\label{sec:cc-rrt-mech}

设定线性系统 $\mathbf{x}_{k+1}=\mathbf{A}\mathbf{x}_k+\mathbf{B}\mathbf{u}_k+\mathbf{w}_k$，$\mathbf{w}_k\sim\mathcal{N}(\mathbf{0},\boldsymbol{\Sigma}_w)$；障碍用半空间 $a_j^\top\mathbf{x}\le b_j$ 描述。

\begin{description}
  \item[机制一：闭环协方差传播。] RRT 每生长一条边（一段标称轨迹），就沿这条边把状态协方差按 \cref{eq:cc-cov-prop} 往前传。若还有障碍位置不确定性 $\boldsymbol{\Sigma}_o$，检验时用 $\mathbf{P}_k+\boldsymbol{\Sigma}_o$（假设独立）。关键：用闭环（含反馈 $\mathbf{K}$）而非开环传播，否则协方差发散。
  \item[机制二：概率可行性检验。] 节点 $\bar{\mathbf{x}}_k$（均值）相对障碍面安全的概率 $\ge1-\delta_{j,k}$，由 \cref{thm:cc-gauss} 给出可直接检验的确定性条件
  \begin{equation}
    a_j^\top\bar{\mathbf{x}}_k-b_j\ \le\ -\Phi^{-1}(1-\delta_{j,k})\sqrt{a_j^\top(\mathbf{P}_k+\boldsymbol{\Sigma}_o)a_j}.
    \label{eq:cc-rrt-feasible}
  \end{equation}
  生长出的新节点若\emph{不满足}此式（违反概率 $>\delta_{j,k}$），直接剔除，不加入树。这一步把"碰撞检测"从布尔的"点在不在障碍里"升级成概率的"违反概率超不超 $\delta$"。
  \item[机制三：Boole 风险分配。] 整条路径的 JCC 由 \cref{thm:cc-boole} 保证——沿路径把 $\delta_{j,k}$ 分配下去（均分或按段分配），使 $\sum_{j,k}\delta_{j,k}\le\delta$。
\end{description}

\begin{insight}
CC-RRT $=$ RRT $+$ 协方差传播 $+$ 概率剔除。确定性 RRT 问"这个点撞不撞"，CC-RRT 问"这个点的违反概率超不超 $\delta$"——把碰撞检测从"几何包含"换成"概率收紧检验"，其余 RRT 机制（最近邻、steer、到达判定）原样保留。这和 \cref{sec:cc-mpc} 把 MPC 的硬约束换成机会约束收紧，是\emph{同一个替换}用在不同规划范式上：优化型 $\to$ CC-MPC，采样型 $\to$ CC-RRT。
\end{insight}

\paragraph{CC-RRT$^*$：渐近最优 + 风险加权。}
把 CC-RRT 嫁接到 RRT$^*$（\cref{sec:plan-rrtstar}）上，继承 \texttt{ChooseParent} 与 \texttt{Rewire}，获得渐近最优。两个关键改动：(1) \emph{风险加权代价}——边代价从 $\|\mathbf{u}\|^2$ 改为 $\|\mathbf{u}\|^2+\lambda\sum_k\delta_{j,k}$（既最短/最省力、又少冒险）；Luders 等论证该风险目标在 RRT$^*$ 框架内可用，故保留渐近最优性。(2) \emph{Rewire 后重检}——rewire 改变父边、从而改变协方差传播路径，\emph{必须重新做机会约束检验}，否则可能引入违反概率超标的连接。

\begin{codebox}[Python]{CC-RRT 的概率可行性检验与 extend}
import numpy as np
from scipy.stats import norm

def propagate_covariance(P, A_cl, Sigma_w):
    return A_cl @ P @ A_cl.T + Sigma_w     # closed-loop one-step

def cc_feasible(x_mean, P, obstacles, delta_step, Sigma_o=None):
    Pe = P if Sigma_o is None else P + Sigma_o
    z = norm.ppf(1.0 - delta_step)         # analytic Gaussian factor
    for a_j, b_j in obstacles:
        margin = z * np.sqrt(float(a_j @ Pe @ a_j))
        if a_j @ x_mean - b_j > -margin:   # violates tightened constraint -> prune
            return False
    return True

def cc_rrt_extend(tree, x_rand, A_cl, Sigma_w, obstacles, delta_step, dt, Sigma_o=None):
    node  = nearest(tree, x_rand)
    x_new = steer(node.mean, x_rand, dt)
    P_new = propagate_covariance(node.P, A_cl, Sigma_w)
    if cc_feasible(x_new, P_new, obstacles, delta_step, Sigma_o):
        return add_node(tree, x_new, P_new, parent=node)
    return None
\end{codebox}

\subsection{协方差控制：不只规划均值，还塑形不确定性}
\label{sec:cc-cs}

前面的方法都只规划\emph{均值轨迹}，协方差是被动传播的副产品。\textbf{协方差控制}（Covariance Steering, CS；Okamoto--Goldshtein--Tsiotras\cite{okamoto2018optimal}）更进一步：把协方差 $\boldsymbol{\Sigma}_k$ 也当成可主动塑形的对象——在宽阔区域"放松"协方差（省控制能量），在狭缝处"收缩"协方差（保证安全）。直觉就像开车过窄桥前主动减速收拢、过桥后再放开。控制律取 $\mathbf{u}_k=\bar{\mathbf{u}}_k+\mathbf{K}_k(\mathbf{x}_k-\bar{\mathbf{x}}_k)$——$\bar{\mathbf{u}}_k$ 操控均值，反馈 $\mathbf{K}_k$ 操控协方差。

\begin{theorem}[有机会约束时均值与协方差不可解耦]\label{thm:cc-cs-coupling}
对线性高斯系统，无机会约束时均值子问题（优化 $\bar{\mathbf{u}}_k$）与协方差子问题（优化 $\mathbf{K}_k$）可分开解。但加入解析高斯收紧 \cref{eq:cc-gauss-tighten} 后，$a^\top\bar{\mathbf{x}}_k\le b-\Phi^{-1}(1-\delta)\sqrt{a^\top\boldsymbol{\Sigma}_k a}$ 中均值 $\bar{\mathbf{x}}_k$ 与协方差 $\boldsymbol{\Sigma}_k$ 同时出现在同一约束里——收紧 margin 依赖 $\boldsymbol{\Sigma}_k$、而 $\boldsymbol{\Sigma}_k$ 由 $\mathbf{K}_k$ 决定。故二者\emph{耦合}，必须\emph{联合求解}。
\end{theorem}

这个不可解耦正是 CS 比"先规划均值、再被动传协方差"更强的原因：它允许"为了让均值能过窄缝，主动把协方差收小"。CS 问题可转化为一个半定规划（SDP），目标含控制反馈能量 $\operatorname{tr}(\mathbf{R}\mathbf{K}_k\boldsymbol{\Sigma}_k\mathbf{K}_k^\top)$ 与状态散布 $\operatorname{tr}(\mathbf{Q}\boldsymbol{\Sigma}_k)$，约束含协方差/均值动力学、机会约束及终端协方差 $\boldsymbol{\Sigma}_N\preceq\boldsymbol{\Sigma}_f$；经变量代换（把 $\mathbf{K}_k\boldsymbol{\Sigma}_k$ 整体作为新变量、Schur 补处理二次项）化成 SDP 标准形\cite{okamoto2018optimal}。后续 Pilipovsky--Tsiotras 用 IRA 做最优风险分配进一步改进。

\begin{insight}
CS 把"不确定性"变成了控制量。前面的机会约束方法把协方差当成已知、被动传播的量，只规划均值绕开它撑大的 margin；CS 反过来主动控制协方差的形状，让不确定性"该小则小"。这是从"适应不确定性"到"塑造不确定性"的范式跃迁。代价是从 SOCP（只规划均值）升级到 SDP（联合规划均值 + 协方差），计算更重；收益是在狭窄/高风险场景能找到解耦方法找不到的可行解。
\end{insight}

\subsection{机会约束 vs 鲁棒 tube：同一收紧思想的两种风险姿态}
\label{sec:cc-vs-tube}

把鲁棒 tube MPC 与机会约束 MPC 放在同一问题上对比，最能看清两者关系。两者都是"名义轨迹 + 约束收紧"，差别只在收紧量怎么来：tube 用"覆盖最坏有界扰动"的不变集半宽 $\Omega$（固定常数），机会约束用"覆盖 $1-\delta$ 概率质量"的 $\Phi^{-1}(1-\delta)\sqrt{a^\top\mathbf{P}_k a}$（随 $\mathbf{P}_k$、$\delta$ 时变）。

\begin{note}
\textbf{tube 是机会约束的极限.}\;机会约束在 $\delta\to0$ 时收紧量 $\Phi^{-1}(1-\delta)\to\infty$——趋于无穷保守，对应"高斯尾部要 100\% 覆盖不可能"。而若把高斯\emph{截断}成有界集，机会约束在 $\delta\to0$ 就退化回 tube。所以\emph{tube 是机会约束在"有界扰动 + $\delta\to0$"下的极限}，机会约束是 tube 在"无界扰动 + 容忍 $\delta$ 风险"下的推广。一维走查：约束 $p\le1$，有界扰动 $|w|\le0.3$、收缩 $\rho=0.4$ 给 tube 半宽 $\Omega=0.3/(1-0.4)=0.5$（恒收 $0.5$）；高斯 $\sigma_\infty=0.2$ 时机会约束 $\delta=0.05$ 收 $0.329<0.5$（更进取），但 $\delta=0.001$ 收 $0.618>0.5$（反更保守）。存在 $\delta^\star$ 使两者相等，两侧各有所长。\pz{控制部深讲：tube MPC 的 RPI 集 $\Omega$ 构造、homothetic/rigid tube、与本节收紧的精确对应，待 \cref{ch:control} 鲁棒 MPC 章；此处只给规划用途的直觉与极限关系。}
\end{note}

\paragraph{前沿落地。}
机会约束近年从仿真走向实机：四足摩擦锥机会约束化（Trivedi et al.\cite{trivedi2025friction}）把单刚体动力学的负载/地形不确定性建模为分布、摩擦锥约束写成机会约束，用高效 QP 求解，Unitree Go1 实机扛起超过自重 50\% 的未知负载、无需重新调参；CC-VPSTO 把机会约束嵌入途经点随机轨迹优化落到机械臂；Robust-RRT 把 CC-RRT 的线性高斯假设放松到不确定非线性系统。\pz{核对结论：四足机会约束论文（Trivedi et al., arXiv:2411.03481）的实机数值已核实——Unitree Go1 实机扛起约 7.3\,kg（机体 $\approx$12\,kg，逾自重 50\%）、跨室内外地形且"无需额外调参"，与正文一致。但出处需更正：该文为 arXiv 预印本、投稿 RA-L 审稿中，\emph{非}已发表的"RA-L 2025"（bib 已按预印本著录）。CC-VPSTO、Robust-RRT 的确切年份/出处仍待核。}

\begin{practice}
\begin{enumerate}
  \item 实现最小 CC-RRT（闭环协方差传播 + Boole 风险分配 + 概率剔除），在 2D 窄通道验证找到的路径蒙特卡洛违反率 $\le\delta$。
  \item 在上面基础上加 RRT$^*$ 的 ChooseParent/Rewire（风险加权代价 + rewire 后重检），对比 CC-RRT 与 CC-RRT$^*$ 的路径代价随采样数的收敛。
  \item 用 SDP 求解器实现一个 1D/2D 狭缝穿越的协方差控制：初始大协方差、缝口收缩、过缝放开，可视化误差椭球沿轨迹的形状变化；说明为什么有机会约束时均值与协方差不可解耦（\cref{thm:cc-cs-coupling}）。
\end{enumerate}
\end{practice}

% ════════════════════════════════════════════════════════════════
\section{POMDP 形式化与信念空间}
\label{sec:pomdp-belief}

\paragraph{动机：状态看不清，怎么办。}
前面的机会约束默认状态 $\mathbf{x}$ \emph{可观测}——只是其演化带概率扰动。但有一类不确定性更根本：\emph{连你自己的状态都看不清}。机器人在长得一样的走廊里定位，传感器读数无法区分"在 A 段"还是"在 B 段"；抓取时手指挡住物体，不知道物体确切位姿；经典的 Tiger 问题\cite{kaelbling1998planning}：两扇门一扇后是老虎、一扇后是奖励，你不知道哪扇，可以"听"获取信息但听不准、且听有代价——该开哪扇门？要不要先多听几次？这些问题的共同点：\textbf{状态不可直接观测，只能通过带噪声的观测去推断}。这正是 POMDP（Partially Observable Markov Decision Process）要处理的。

\paragraph{反面：把 POMDP 当 MDP 做会怎样。}
最朴素的偷懒：用滤波器（EKF/粒子滤波）估一个点估计 $\hat s$（均值或 MAP），假装它是真实状态、套 MDP 最优策略。两个典型失败：\emph{失败一}——走廊两端一样时，belief 可能\emph{双峰}，点估计（均值）落在两峰中间一个根本不可能的位置（走廊正中的墙里），决策荒谬。\emph{失败二}——MDP 策略\emph{没有"为了看得更清而行动"的概念}（它假设状态本就可观测），于是 Tiger 里当 MDP 不确定时，要么瞎猜开门、要么僵局，而不会"先多听几次"——"听"在 MDP 里是无收益动作（不改变物理状态、只改变信息），MDP 看不到它的价值。

\begin{insight}
POMDP 的关键是"在分布上决策"，而非"在点估计上决策"。把 belief 塌缩成一个点再做 MDP，丢掉的恰恰是"我有多不确定"这个信息——而它正是 POMDP 决策的核心依据。不确定性大时，最优 POMDP 策略倾向于\emph{先降低不确定性}（获取信息）再行动；不确定性小时才直接追求奖励。这种"该探索还是该利用、由当前 belief 的形状决定"的行为，是 MDP 永远学不会的，因为 MDP 的世界里没有"不确定的状态"这回事。正是这一点让 POMDP 能自然导出主动感知 / Active SLAM（\cref{sec:pomdp-bsp}）。
\end{insight}

\subsection{POMDP 七元组与 belief 的贝叶斯滤波}
\label{sec:pomdp-seven}

\begin{definition}[POMDP 七元组]\label{def:pomdp-seven}
POMDP 在 MDP 五元组 $\langle S,A,T,R,\gamma\rangle$ 上加两项（观测空间和观测模型），成为七元组 $\langle S,A,T,R,\Omega,O,\gamma\rangle$：状态空间 $S$（\emph{往往不可直接观测}）、动作空间 $A$、转移 $T(s'\mid s,a)$、奖励 $R(s,a)$、观测空间 $\Omega$、观测模型 $O(o\mid s',a)$（在 $a$ 后到达 $s'$ 时观测到 $o$ 的概率）、折扣因子 $\gamma\in[0,1)$。
\end{definition}

与 MDP 唯一的结构差别是 $\Omega$ 和 $O$：你不再直接看到 $s'$，而是看到由 $s'$ 概率性生成的观测 $o$。$O(o\mid s',a)$ 编码传感器的不完美——理想传感器 $O$ 确定（$o$ 唯一确定 $s'$），现实传感器 $O$ 弥散（同一 $s'$ 可能给出不同 $o$，不同 $s'$ 可能给出相同 $o$，即感知混淆）。

既然看不到状态，就维护一个对状态的概率分布——\textbf{信念}（belief）$b$，$b(s)$ 是"当前处于状态 $s$ 的概率"，$\sum_s b(s)=1$。每走一步（执行动作 $a$、收到观测 $o$），belief 按\textbf{贝叶斯滤波}更新：
\begin{equation}
  b'(s')=\eta\,\underbrace{O(o\mid s',a)}_{\text{观测校正}}\,\underbrace{\sum_{s}T(s'\mid s,a)\,b(s)}_{\text{运动预测}},
  \label{eq:pomdp-belief-update}
\end{equation}
其中 $\eta=1/\Pr(o\mid b,a)$ 是归一化常数。逐项：\emph{运动预测} $\sum_s T(s'\mid s,a)b(s)$ 把旧 belief 通过转移模型推一步（\emph{增大}不确定性）；\emph{观测校正} $O(o\mid s',a)\cdot$ 用观测 $o$ 给预测 belief 重新加权（\emph{减小}不确定性）；\emph{归一化} $\eta$ 让结果重新成为概率分布。

\begin{note}
\textbf{belief update 就是贝叶斯滤波.}\;它统一了你可能已熟悉的几个框架：HMM 前向算法是它的"离散无动作"版；卡尔曼滤波（\cref{sec:ctrl-lqe-dual}）是它在"线性高斯"下的解析特例（belief 是高斯 $(\boldsymbol{\mu},\boldsymbol{\Sigma})$、预测/更新有闭式 Riccati 递归）；粒子滤波是它的非参数蒙特卡洛实现。EKF 的"预测步"$=$ 运动预测、"更新步"$=$ 观测校正。\emph{如果你写过 EKF-SLAM，你已经实现过 belief update}——区别只在 SLAM 到此为止（估完用均值），POMDP 还要在 $b'$ 之上做规划。
\end{note}

\begin{pitfall}[编程陷阱：belief update 忘了归一化，或归一化在校正前]
现象：belief 不再是概率分布（和不为 1），多步后数值漂移甚至下溢全为 0。错误做法：算完 $O\cdot\sum T\cdot b$ 直接用（没除 $\eta^{-1}$），或在乘 $O$ 之前就归一化。正确写法：先算未归一化的 $\tilde b(s')=O(o\mid s',a)\sum_s T(s'\mid s,a)b(s)$，\emph{全部算完后}再统一除以 $\sum_{s'}\tilde b(s')$——因为 $\eta$ 依赖所有 $s'$ 的未归一化值。检验：每步 update 后断言 $|\sum_s b(s)-1|<10^{-9}$；大状态空间用对数域防下溢。
\end{pitfall}

\subsection{belief 是充分统计量，POMDP = belief 空间的 MDP}
\label{sec:pomdp-suff}

这是 POMDP 理论的枢纽。

\begin{theorem}[belief 是充分统计量，POMDP 化归 belief-MDP]\label{thm:pomdp-suff}
belief 是历史的\textbf{充分统计量}：给定当前 belief $b$，过去整段动作-观测历史 $(a_1,o_1,\dots,a_t,o_t)$ 可\emph{完全遗忘}——未来的最优决策只依赖 $b$，不依赖到达 $b$ 的历史。两段不同历史若导出相同 belief，之后的最优策略完全一样。由此 POMDP 可转化为 belief 空间上的（完全可观测）MDP——belief-MDP：状态 $=$ belief $b$（\emph{完全可观测}，因为它是你自己算出来的），动作 $=$ 原动作 $a$，转移 $=$ belief update（$b'=\tau(b,a,o)$，概率 $\Pr(o\mid b,a)$），奖励 $\rho(b,a)=\sum_s b(s)R(s,a)$。其 Bellman 最优方程为
\begin{equation}
  V^*(b)=\max_{a}\Big[\rho(b,a)+\gamma\sum_{o}\Pr(o\mid b,a)\,V^*\big(\tau(b,a,o)\big)\Big].
  \label{eq:pomdp-bellman}
\end{equation}
\end{theorem}
\begin{proof}[证明思路]
$b$ 是状态的完整后验，已把历史里所有"与未来相关"的信息编码进去（由 \cref{eq:pomdp-belief-update} 的递归性，$b'$ 只依赖 $b,a,o$ 而不依赖更早历史）。故对任意策略，未来期望回报仅是 $b$ 的函数，Bellman 最优方程 \cref{eq:pomdp-bellman} 直接套用 \cref{thm:plan-bellman} 的最优性原理于 belief-MDP 即得。
\end{proof}

\begin{insight}
POMDP $\to$ belief-MDP 是"用状态空间的复杂度换可观测性"。原 POMDP 的麻烦是"状态不可观测"；转成 belief-MDP 后状态（belief）完全可观测了——天大的好消息。但代价是状态空间从离散有限集 $S$（$|S|$ 个状态）变成 belief 所在的\emph{连续高维单纯形} $\Delta^{|S|-1}$。你把"看不清"的难题，转化成了"在一个连续高维空间上做 MDP"的难题。POMDP 的全部求解技术（\cref{sec:pomdp-alpha,sec:pomdp-solve}）都在跟这个连续高维 belief 空间搏斗。
\end{insight}

\begin{pitfall}[思维陷阱：以为"观测越多 belief 一定越确定"]
"多观测几次，不确定性总会下降吧"——错。在感知混淆环境（走廊两端一样）里反复观测，belief 始终双峰、不收敛。根因：观测能降低不确定性的前提是它对不同状态\emph{有区分力}（$O(o\mid s_1)\ne O(o\mid s_2)$）；感知混淆区观测对两状态似然相同，再多也不消歧——必须靠\emph{动作}（移动到能区分的地方）才能消歧。这正是 POMDP / Active SLAM 要解决的。
\end{pitfall}

\begin{practice}
\begin{enumerate}
  \item 在 Tiger 问题上实现 belief update：从均匀 belief $(0.5,0.5)$ 出发，连续"听"到 3 次同侧观测（如 $O(\text{hear-left}\mid\text{tiger-left},\text{listen})=0.85$），打印每步 belief，验证它向该侧收敛但不到 1。
  \item 构造一个 4 状态走廊，其中状态 1 和 3 观测模型完全相同（感知混淆）。从在 1、3 上各半的 belief 出发反复"原地观测"，验证 belief 不收敛（双峰）；再加"前进"动作离开混淆区，验证这时才收敛。
  \item 论证"两段不同历史若导出相同 belief，则之后最优策略相同"，并说明它为什么是"POMDP 可转化为 belief-MDP"的前提。
\end{enumerate}
\end{practice}

% ════════════════════════════════════════════════════════════════
\section{\texorpdfstring{$\alpha$}{alpha}-vector、PWLC 值函数与两类灾难}
\label{sec:pomdp-alpha}

\paragraph{动机：belief 空间上的值函数长什么样。}
\cref{thm:pomdp-suff} 把 POMDP 化成 belief 空间上的 MDP，Bellman 方程也照搬了。但 belief 空间是\emph{连续高维}的（$|S|$ 维单纯形），值函数 $V^*(b)$ 定义在这个连续空间上——它长什么样？能否用有限的东西表示？若不能，POMDP 就根本无从求解。Sondik 在 1971 年的博士论文\cite{sondik1971optimal} 回答了这个问题。

\begin{theorem}[Sondik：有限视界值函数 PWLC]\label{thm:pomdp-pwlc}
有限视界 POMDP 的最优值函数 $V^*(b)$ 在 belief 空间上是\textbf{分段线性且凸的}（Piecewise-Linear and Convex, PWLC），可写成有限个线性函数取上包络：
\begin{equation}
  V^*(b)=\max_{\boldsymbol{\alpha}\in\Gamma}\ \sum_{s}\alpha(s)\,b(s)=\max_{\boldsymbol{\alpha}\in\Gamma}\ \boldsymbol{\alpha}\cdot b,
  \label{eq:pomdp-alpha}
\end{equation}
其中 $\Gamma=\{\boldsymbol{\alpha}_1,\dots,\boldsymbol{\alpha}_m\}$ 是有限个 $\boldsymbol{\alpha}$-vector 的集合，每个 $\boldsymbol{\alpha}$ 是一个 $|S|$ 维向量。
\end{theorem}

逐项理解：每个 $\boldsymbol{\alpha}$-vector 是 belief 空间上的一个\emph{线性函数} $b\mapsto\boldsymbol{\alpha}\cdot b$（单纯形上的一张超平面）；$V^*(b)$ 取所有这些线性函数的逐点最大（上包络）——有限个线性函数取 max，结果分段线性且凸；在每个 belief 区域，某一个 $\boldsymbol{\alpha}$-vector 占主导。\emph{为什么凸}（直觉）：值函数凸意味着"belief 越不确定（越靠单纯形中心），价值越低"——不确定时更可能选错动作；在顶点（belief 集中在某确定状态）价值最高。

\paragraph{$\boldsymbol{\alpha}$-vector $=$ 条件计划。}
$\boldsymbol{\alpha}$-vector 有清晰语义：\emph{每个 $\boldsymbol{\alpha}$-vector 对应一个条件计划}（conditional plan）——"从某个动作开始、之后根据观测如何分支"的一整套策略；$\alpha(s)$ 的值 $=$ "如果真实状态是 $s$，执行这个条件计划能拿到的期望累积奖励"。于是最优策略极简：在当前 belief $b$ 下，选使 $\boldsymbol{\alpha}\cdot b$ 最大的那个 $\boldsymbol{\alpha}$-vector、执行它对应的动作——$\Gamma$ 同时编码值函数（取上包络）和策略（取 argmax 的那个 $\boldsymbol{\alpha}$ 对应的动作）。

\begin{note}
\textbf{2 状态走查.}\;Tiger 的 belief 是一个数 $p=b(\text{TL})\in[0,1]$，每个 $\boldsymbol{\alpha}$-vector 是 $[0,1]$ 上一条直线。设三个条件计划：开右门 $\boldsymbol{\alpha}_A=(10,-100)$、开左门 $\boldsymbol{\alpha}_B=(-100,10)$、先听 $\boldsymbol{\alpha}_C=(8.5,8.5)$（近水平线）。$V^*(p)=\max$ 这三条直线的上包络分三段：$p$ 极小开左、中段听、$p$ 极大开右（如 $\boldsymbol{\alpha}_A$ 与 $\boldsymbol{\alpha}_C$ 交于 $10p-100(1-p)=8.5\Rightarrow p\approx0.986$）——一条分段线性、凸的曲线，每段由一个条件计划主导。这就是"$p\approx0.99$ 才开门、否则听"的几何来由。\pz{存疑：这些 $\boldsymbol{\alpha}$-vector 数值（$10,-100,8.5$ 等）与交点 $0.986$ 是源 md 为教学构造的示意数，非 Tiger 标准参数下的精确解；作教学走查可，勿当权威数值。}
\end{note}

\subsection{两类灾难：为什么精确求解不可扩展}
\label{sec:pomdp-curse}

既然 $V^*$ 能用有限个 $\boldsymbol{\alpha}$-vector 精确表示，为什么 POMDP 仍出名地难？因为 $\boldsymbol{\alpha}$-vector 的数量随视界指数爆炸。

\begin{insight}
POMDP 的难，难在"两类 curse 叠加"。\textbf{维数灾难（curse of dimensionality）}：belief 是 $|S|$ 维连续向量，状态一多 belief 空间就爆炸。\textbf{历史灾难（curse of history）}：做一步 Bellman 备份时，新 $\boldsymbol{\alpha}$-vector 集合要枚举"当前动作 $\times$ 每个观测分别接哪个旧 $\boldsymbol{\alpha}$-vector"的所有组合，规模以 $|\Gamma|^{|\Omega|}$ 的方式膨胀（观测分支带来）。精确求解同时撞上这两堵墙（被证明 PSPACE-hard 量级）。后续所有近似方法本质都是在绕开其中一堵或两堵：点基方法（\cref{sec:pomdp-pbvi}）绕维数灾难（只在少数可达 belief 点上备份）；在线树搜索 DESPOT（\cref{sec:pomdp-despot}）同时绕两堵（采样状态绕维数灾难、采样观测绕历史灾难）。
\end{insight}

实践中绝大多数新生成的 $\boldsymbol{\alpha}$-vector 是\emph{冗余的}（在任何 belief 上都不是上包络的一部分，被其他 $\boldsymbol{\alpha}$ 支配）。精确算法（Sondik one-pass、Witness、Incremental Pruning）大量时间花在"剪掉被支配的 $\boldsymbol{\alpha}$-vector"上，但即便如此也扛不住指数增长——几十个状态、十几步视界就难以为继。

\paragraph{QMDP：最廉价的近似。}
理解了 $\boldsymbol{\alpha}$-vector，一个极廉价的近似立刻浮现：\textbf{QMDP}。它假设"再走一步后世界就完全可观测了"——用底层 MDP 的 $Q$ 函数 $Q_{\mathrm{MDP}}(s,a)$ 当 $\boldsymbol{\alpha}$-vector：$\alpha_a(s)=Q_{\mathrm{MDP}}(s,a)$，决策 $a^*=\arg\max_a\sum_s b(s)Q_{\mathrm{MDP}}(s,a)$。QMDP 极快（只需解一个 MDP），但有系统性短板：\emph{它系统性低估"持续获取信息"的长期价值}。因为它假设"再走一步世界就全可观测"，于是消歧不足、倾向过早行动。需澄清一个常见误传：QMDP 并非"从不获取信息"（某些 belief 上它仍会选信息动作），它真正的毛病是把"消歧要做多久"压缩成"一步就够"。结论：QMDP 适合"信息会自然到来、只需少量消歧"的问题（快且够用），不适合需要大量主动消歧的问题。

\begin{pitfall}[思维陷阱：以为 $\boldsymbol{\alpha}$-vector 爆炸是"实现不好"，换算法就行]
"精确求解慢是因为算法没优化，用更好的剪枝就能扩展"——错。$\boldsymbol{\alpha}$-vector 随视界指数增长是 POMDP 的本质难度（两个 curse），再好的剪枝也只能延缓、不能消除指数（PSPACE-hard 是问题固有的）。正确做法：中等规模上点基近似（SARSOP，\cref{sec:pomdp-pbvi}）；大规模/实时上在线树搜索（DESPOT，\cref{sec:pomdp-despot}）。精确求解只用于很小的问题或理论分析。
\end{pitfall}

\begin{pitfall}[概念误区：把 $\boldsymbol{\alpha}$-vector 当 belief（同为 $|S|$ 维向量，易混）]
belief $b$ 是\emph{概率分布}（非负、和为 1），表示"我在哪"；$\boldsymbol{\alpha}$-vector 是\emph{价值向量}（可正可负、不归一化），$\alpha(s)$ 表示"若真在 $s$、执行某计划的价值"。两者维度相同但语义完全不同，通过内积 $\boldsymbol{\alpha}\cdot b$ 配对。检验：$\boldsymbol{\alpha}$ 的分量出现负数是正常的（奖励可负），belief 出现负数是 bug。
\end{pitfall}

\begin{practice}
\begin{enumerate}
  \item 在 Tiger 上实现 QMDP：先解全可观测 MDP 得 $Q_{\mathrm{MDP}}(s,a)$，再用 $a^*=\arg\max_a\sum_s b(s)Q_{\mathrm{MDP}}(s,a)$ 决策。观察它在不同 belief 下何时 listen、何时开门，理解它系统性低估"持续消歧"价值。
  \item 取一个 2 状态 POMDP，给定 3 个 $\boldsymbol{\alpha}$-vector（各是一条直线），画出上包络 $V^*(p)=\max_{\boldsymbol{\alpha}}\boldsymbol{\alpha}\cdot b$，标出每段由哪个 $\boldsymbol{\alpha}$ 主导、对应哪个动作，验证上包络是凸的分段线性曲线。
  \item 为什么 $V^*(b)$ 是\emph{凸}而非凹？从"belief 越不确定、价值越低"的直觉和"有限个线性函数取 max"的代数两方面论证。
\end{enumerate}
\end{practice}

% ════════════════════════════════════════════════════════════════
\section{从精确到近似：点基方法与在线树搜索}
\label{sec:pomdp-solve}

\subsection{点基方法：只在可达 belief 上备份}
\label{sec:pomdp-pbvi}

\paragraph{动机：不必在整个单纯形上求解。}
精确求解的死穴是"在整个 belief 单纯形上维护 $\boldsymbol{\alpha}$-vector"。但有一个救命观察：\emph{从固定初始 belief $b_0$ 出发，实际可达的 belief 只是单纯形里的一小块}——由"从 $b_0$ 出发、各种动作-观测序列能到达的 belief"构成的低维子集。\textbf{点基方法}（point-based methods）正是抓住这一点：只在一组\emph{代表性的可达 belief 点}上做 Bellman 备份，每个点维护一个 $\boldsymbol{\alpha}$-vector。

\cref{tab:pomdp-pbvi} 给出三代点基方法的演进——核心差别在"往哪放 belief 点 / 往哪备份"。

\begin{table}[H]\centering\small
\caption{点基价值迭代三代演进}
\label{tab:pomdp-pbvi}
\begin{tabular}{lll}
\toprule
方法 & 选 belief 点的策略 & 关键贡献 \\
\midrule
PBVI\cite{pineau2003pbvi} & 从 $b_0$ 扩展 belief 集、覆盖可达区域 & 开创点基；$\boldsymbol{\alpha}$ 数 $=$ 点数（线性，绕维数灾难）\\
HSVI\cite{smith2004hsvi} & 沿"上下界 gap 最大"路径下探 & 上下界引导 + 有界次优保证 \\
SARSOP\cite{kurniawati2008sarsop} & 只在"最优可达"belief 上备份 & 离线 $\epsilon$-最优事实标准求解器 \\
\bottomrule
\end{tabular}
\end{table}

三者层层聚焦：PBVI"覆盖可达"$\to$ HSVI"优先 gap 大的可达"$\to$ SARSOP"只盯最优策略会去的可达"。越往后，同样预算花在越值得的地方，逼近越快。

\begin{insight}
点基方法用"可达性"砍掉维数灾难的大头。精确求解的浪费在于"为永远到不了的 belief 也求解"。点基只在从 $b_0$ 可达（SARSOP 更进一步：最优可达）的 belief 上备份——可达 belief 集通常是单纯形里极低维的子集，于是 $\boldsymbol{\alpha}$-vector 数量从"随视界指数"降到"随采样点数线性"。每个 $\boldsymbol{\alpha}$ 对应一个\emph{真实可执行}的条件计划，故 $\tilde V(b)=\max_{\boldsymbol{\alpha}}\boldsymbol{\alpha}\cdot b\le V^*(b)$ 是 $V^*$ 的\emph{下界}（这让点基解"安全"、不会高估价值），且随备份单调上升、逼近 $V^*$。
\end{insight}

\paragraph{离线 vs 在线定位。}
点基方法是\emph{离线}的：预先算出整个策略（一个 $\boldsymbol{\alpha}$-vector 集合），在线执行只需查表（每步取 $\arg\max$ 一次内积，极快）。适用边界：\emph{适合}状态空间中等（$|S|$ 几十到几千）、模型固定、需反复快速决策；\emph{不适合}状态空间巨大、模型在线变化（如自驾每帧场景不同）、初始 belief 经常变（SARSOP 的可达空间针对特定 $b_0$）。这些情况要用\emph{在线方法}（\cref{sec:pomdp-despot}）。离线与在线常\emph{合作}：点基 $\boldsymbol{\alpha}$-vector 下界可当在线搜索的初始下界/rollout 策略（"用粗地图加速精搜"）。

\subsection{在线树搜索：POMCP 与 DESPOT}
\label{sec:pomdp-despot}

\paragraph{动机：当离线行不通。}
状态空间巨大、模型逐帧变化、只需"此刻"的决策时，\textbf{在线 POMDP 规划}登场：不预先算整个策略，而是在每个决策周期\emph{只为当前 belief $b_0$ 现场规划}——从 $b_0$ 出发构造有限深度的搜索树、找出当前最优动作、执行一步，然后下一周期用更新后的 belief 重新规划。

\paragraph{前奏：POMCP。}
在 DESPOT 之前，\textbf{POMCP}（Silver--Veness\cite{silver2010pomcp}）把蒙特卡洛树搜索（MCTS）搬到 belief 树上：用一组状态粒子近似 belief（让它能处理巨大/连续状态，粒子数与 $|S|$ 无关）；MCTS 四步（用 UCT/UCB 选动作平衡探索-利用、采样观测下探、随机 rollout 估叶子、反向传播）。短板：UCB 启发式可能被少数高方差 rollout 带偏，最坏情况差。DESPOT 保留"用采样破两个 curse"的优点，但用一组\emph{固定的采样场景}来评估策略，避免 POMCP 的极端最坏行为。

\begin{definition}[确定化场景与 DESPOT]\label{def:pomdp-despot}
一个\textbf{场景}（determinized scenario）$\phi=(s_0,\phi_1,\phi_2,\dots)$ 是"一个起始状态 + 一串预先采样好的随机数"。给定 $\phi$ 和一个动作序列，状态和观测的演化就\emph{完全确定}（随机数把转移和观测在这次模拟里钉死）。\textbf{DESPOT}（Determinized Sparse Partially Observable Tree；Somani--Ye--Hsu--Lee\cite{somani2013despot}）是标准 belief 树在 $K$ 个采样场景 $\{\phi_1,\dots,\phi_K\}$ 下的稀疏子树：\emph{保留所有动作分支}（每个节点仍可选任意动作），但\emph{只保留 $K$ 个场景实际经历的观测分支}（而非所有 $|\Omega|$ 种观测）。
\end{definition}

完整 belief 树（深度 $D$）节点数 $O(|A|^D|\Omega|^D)$（双指数，历史灾难）。DESPOT 把它降到 $O(|A|^D K)$——观测分支被场景数 $K$ 限制。破两个 curse 的方式：\emph{采样状态}（从 belief 采 $K$ 个起始状态）破维数灾难，\emph{采样观测}（只保留场景经历的观测分支）破历史灾难。评估一个策略 $\pi$，就是看它在 $K$ 个场景上跑出的平均折扣回报 $\hat V_\pi(b)=\frac1K\sum_k V_{\pi,\phi_k}$。

\begin{theorem}[DESPOT 的 output-sensitive regret 界（直觉陈述）]\label{thm:pomdp-output}
从 DESPOT 得到的最优策略是近最优的，其 regret（与真最优的差距）界\emph{依赖于"最优策略的表示大小"，而非状态空间大小}\cite{somani2013despot}。直白说：\emph{只要存在一个简洁（表示小）的好策略，DESPOT 用不多的场景 $K$ 就能找到它}——哪怕状态空间天文数字。这是 DESPOT 能在巨大状态空间实时工作的理论根基。
\end{theorem}

\paragraph{过拟合与 R-DESPOT。}
只看"在 $K$ 个场景上的平均回报"有个陷阱：复杂策略可能在这 $K$ 个特定场景上完美（过拟合），却在真实分布上泛化差。\textbf{R-DESPOT}（Regularized DESPOT）优化正则化目标
\begin{equation}
  \max_{\pi}\ \Big[\hat V_\pi(b_0)-\lambda\,|\pi|\Big],
  \label{eq:pomdp-rdespot}
\end{equation}
其中 $|\pi|$ 是策略（子树）大小，$\lambda$ 是正则化系数——惩罚过大的策略、逼它简洁从而泛化好。

\begin{insight}
R-DESPOT 的正则化 $=$ 机器学习的"奥卡姆剃刀"搬到规划里：$K$ 个场景是"训练集"、策略是"模型"，在 $K$ 个场景上完美但真实分布差就是过拟合，$\lambda|\pi|$ 就是复杂度惩罚。这不是巧合——"用有限样本估计、防过拟合要正则化"是统计学习的普适原理。$K$ 太小时一定要正则化（样本少、过拟合风险高），$\lambda$ 随样本越少越大。
\end{insight}

\cref{tab:pomdp-online} 横向对比几个在线 POMDP 求解器，便于选型（它们共享"在线、从当前 belief 现算、anytime"的框架，分野在"如何花有限的仿真预算"）。

\begin{table}[H]\centering\small
\caption{在线 POMDP 求解器横向对比}
\label{tab:pomdp-online}
\begin{tabular}{lll}
\toprule
求解器 & 核心机制 & 取舍 \\
\midrule
POMCP\cite{silver2010pomcp} & belief 树 MCTS + UCT + 随机 rollout & 简单、大 POMDP 实践好；最坏情况差 \\
DESPOT\cite{somani2013despot} & $K$ 确定化场景稀疏树 + 上下界 + 正则化 & 方差低、output-sensitive；需调 $K/\lambda$ \\
AEMS & 启发式搜索 belief 树（上下界 gap 选下探）& 有界次优、收敛快；大状态吃力 \\
POMCPOW & POMCP + Progressive Widening & 处理连续动作/观测；调参更敏感 \\
\bottomrule
\end{tabular}
\end{table}

\paragraph{DESPOT 家族演化。}
NUS AdaComp 组围绕 DESPOT 做了一条"从能跑通到上自动驾驶车"的演化链\cite{somani2013despot}：HyP-DESPOT（CPU 多线程 + GPU CUDA 并行 rollout，加速百倍，让 DESPOT 进实时控制回路——根本原因是"$K$ 个场景彼此独立"这一设计本就并行友好）；DESPOT-$\alpha$（$\boldsymbol{\alpha}$-vector 粒子近似 + 子策略共享，治大观测空间下的过度乐观）；Context-POMDP（注入道路上下文 + 行人隐藏意图，自驾实例化）；MAGIC（学宏动作砍视界）/ LeTS-Drive（离线模仿树搜索训练网络、在线引导 HyP-DESPOT）。三条轴（硬件/领域/学习）正交、可叠加。\pz{存疑：DESPOT 家族各成员的会议/年份（HyP-DESPOT RSS 2018、Context-POMDP ICRA 2020、MAGIC/LeTS-Drive RSS 2019/2021 等）来自源 md，作脉络叙述可，确切归属待联网核。}

\begin{insight}
DESPOT $=$ "用 $K$ 个确定化场景，把双指数的随机 belief 树压成线性于 $K$ 的稀疏确定树"。完整 belief 树的爆炸来自两处：状态多、观测序列多。DESPOT 用同一招"采样"同时治两病——采 $K$ 个起始状态治前者、只留这 $K$ 个场景经历的观测分支治后者；再用 output-sensitive 界（\cref{thm:pomdp-output}）保证"只要好策略简洁，少量场景就够"，用 R-DESPOT 正则化防过拟合。这让 POMDP 第一次在"$10^5$ 量级状态 + 实时"的真实机器人问题上跑通，成了在线 POMDP 的事实标准。
\end{insight}

\begin{pitfall}[概念误区：以为"场景"是采样的状态序列]
"$K$ 个场景就是 $K$ 条采样出来的状态轨迹吧"——错。场景 $\phi=(s_0,\phi_1,\phi_2,\dots)$ 是"起始状态 + 一串\emph{随机数}"，不是轨迹。随机数把随机性钉死，但走出什么轨迹仍取决于你\emph{选什么动作}——同一场景 + 不同动作序列 $=$ 不同确定轨迹。这正是"一棵 DESPOT 编码所有策略在 $K$ 场景下的执行"的关键。把场景想成"预先掷好的骰子序列"——决策还是你做的，骰子只是不再随机。
\end{pitfall}

\begin{practice}
\begin{enumerate}
  \item 实现"场景 $=$ 起始状态 + 随机数串"的 \texttt{step(state, action, rand\_seq)}：验证同一随机数串 + 不同动作序列给出不同轨迹，理解场景为何能"确定化"随机执行。
  \item 在 Tiger（或 RockSample）上跑一个 DESPOT 式在线规划器，分别用 $K=10,50,200$ 场景，记录策略质量与单步规划耗时，验证 $K$ 越大越好但越慢，并观察 $K$ 太小时的过拟合。
  \item 给一个 $|A|=2,|\Omega|=2,D=3$ 的小 POMDP，分别画出完整 belief 树和 $K=3$ 场景下的 DESPOT，数节点数，直观感受"稀疏"省了多少。
\end{enumerate}
\end{practice}

% ════════════════════════════════════════════════════════════════
\section{信念空间运动规划与 Active SLAM}
\label{sec:pomdp-bsp}

\paragraph{动机：连续高维状态怎么办。}
前面的 POMDP 方法主要面向\emph{离散}状态/观测或用粒子处理连续状态。但机器人有一大类问题是\emph{连续高维}的（机械臂关节位形、移动机器人连续位姿）。在这些问题上，机器人规划界\emph{独立}发展出\textbf{信念空间运动规划}（belief-space motion planning, BSP），核心假设是 \emph{belief 为高斯}（用均值 + 协方差刻画）。高斯假设的代价与收益：收益是 belief 用 $(\boldsymbol{\mu},\boldsymbol{\Sigma})$ 紧凑表示、可微、能用卡尔曼系机器解析处理高自由度；代价是只能表示\emph{单峰}（处理不了多模态歧义）。所以 BSP 是 POMDP 在"高斯 belief + 连续空间"下的特化：擅长高自由度，不擅长多模态。

\cref{tab:bsp} 梳理 BSP 主线。

\begin{table}[H]\centering\small
\caption{信念空间运动规划方法脉络}
\label{tab:bsp}
\begin{tabular}{lll}
\toprule
方法 & 核心思想 & 关键贡献 \\
\midrule
LQG-MP\cite{vandenberg2011lqgmp} & 对候选路径用 LQG 解析传播协方差 & 执行前评估路径的碰撞概率、可比较 \\
ML 观测假设\cite{platt2010belief} & 假设未来观测 $=$ 最大似然值 & belief 规划化归确定性轨迹优化 \\
RRBT\cite{bry2011rrbt} & RRT$^*$ 进 belief 空间（节点带协方差）& 采样式找"既短又确定"的路径 \\
FIRM\cite{aghamohammadi2014firm} & 节点放 stationary LQG 控制器钉死 belief & 节点 belief 与历史无关，破历史灾难 \\
BSP-iSAM\cite{indelman2015bsp} & 因子图里可微算"未来信息增益" & 接 SLAM 增量平滑、贯通 Active SLAM \\
\bottomrule
\end{tabular}
\end{table}

\paragraph{协方差沿路径传播。}
BSP 的"高斯 belief 演化"就是协方差矩阵 $\boldsymbol{\Sigma}$ 沿路径的卡尔曼递归（\cref{sec:ctrl-lqe-dual}）：预测步 $\boldsymbol{\Sigma}^-=\mathbf{A}\boldsymbol{\Sigma}\mathbf{A}^\top+\mathbf{W}$（不确定性增大），更新步 $\boldsymbol{\Sigma}^+=(\mathbf{I}-\mathbf{K}\mathbf{H})\boldsymbol{\Sigma}^-$（观测带来信息、不确定性减小）。LQG-MP 沿候选路径反复跑这个递归，预先算出每个路点的 $\boldsymbol{\Sigma}$；有了 $\boldsymbol{\Sigma}$，路点的碰撞概率可由"障碍方向上的位置方差"经高斯尾概率估出（距离 $d$、投影标准差 $\sigma$，碰撞概率 $\approx\Phi(-d/\sigma)$）——这与 \cref{thm:cc-gauss} 是同一个高斯尾概率思想。FIRM 的巧招：在每个图节点放 LQG 控制器，让机器人稳定到一个固定协方差 $\boldsymbol{\Sigma}_\infty$（离散 Riccati 不动点），于是"到达节点后的 belief"被钉成固定值、\emph{与到达路径无关}——把"带历史的 belief 传播"退化成普通图搜索，历史灾难被破。

\begin{insight}
BSP 是 POMDP 的"高斯特例"，像卡尔曼滤波是贝叶斯滤波的高斯特例。通用 POMDP 对任意 belief 成立（$\boldsymbol{\alpha}$-vector/粒子，能处理多模态但贵），BSP 是它在"高斯 belief"下的特例（用 $(\boldsymbol{\mu},\boldsymbol{\Sigma})$ 解析处理高自由度，但只能单峰）。所以 BSP 与 DESPOT 不是竞争而是互补——高自由度连续问题用 BSP，多模态/感知混淆问题用 DESPOT。
\end{insight}

\subsection{Active SLAM = 信息增益作奖励的 POMDP}
\label{sec:pomdp-active}

经典 SLAM 是\emph{被动}的：机器人沿给定路径走，顺便建图。但有个更主动的问题：\emph{"下一步该去哪里看，才能最快把地图建准、把自己定准？"} 这就是 \textbf{Active SLAM}。它本质是一个我们已经会的决策问题——POMDP。把 Active SLAM 套进 \cref{def:pomdp-seven} 的七元组：状态 $s$ $=$ 机器人位姿 + 地图特征（连续高维）；动作 $a$ $=$ 去哪里看；观测 $o$ $=$ 传感器读数；转移 $T$、观测 $O$ $=$ 运动/传感器模型；belief $b$ $=$ 位姿 + 地图的联合后验（\emph{就是 SLAM 的输出}）；奖励 $R$ $=$ \textbf{信息增益}（让 belief 更确定）+ 覆盖 $-$ 碰撞 $-$ 运动代价。\emph{唯一新的东西是"奖励 $=$ 信息增益"——其余全是 SLAM 已有的}。

怎么量化"belief 更确定"？用协方差 $\boldsymbol{\Sigma}$（或信息矩阵 $\boldsymbol{\Sigma}^{-1}$）的标量化，即最优实验设计的几个准则：\textbf{D-optimality} $\det(\boldsymbol{\Sigma})$（或 $\log\det$，不确定性椭球\emph{体积}）、\textbf{A-optimality} $\operatorname{tr}(\boldsymbol{\Sigma})$（各方向方差之\emph{和}）、\textbf{E-optimality} $\lambda_{\max}(\boldsymbol{\Sigma})$（椭球\emph{最长轴}）。信息增益 $=$ 行动前后这个标量的下降，也可用 Shannon 熵的下降（互信息）度量\cite{placed2023survey}。

\begin{insight}
Active SLAM 让"SLAM 技能"长出"决策能力"，几乎零额外成本。SLAM 工程师已经有了 POMDP 七元组的六个元素（状态、动作、转移、观测、belief update），唯一缺的是"奖励 $=$ 信息增益"和"在 belief 上规划"。补上这两块，被动 SLAM 就升级成主动 SLAM——机器人自己决定"去哪看最划算"。BSP-iSAM 的"在因子图里可微算未来信息增益"正是把估计与决策无缝缝合的工程实现。
\end{insight}

\paragraph{信息增益为何贵、为何贪心没那么糟。}
信息增益 $\mathrm{IG}(b,a)=H(b)-\mathbb{E}_o[H(\tau(b,a,o))]$ 里有一个对\emph{未来观测求期望}的 $\mathbb{E}_o$——这正是它贵的根源（要枚举/采样动作后可能收到的每个观测、各做一次 belief update、算熵、加权）。多步前瞻更糟（嵌套期望、指数上升）。所以实用系统多用单步贪心近似：next-best-view（贪心选"能看到最多新信息的视角"）、frontier 探索（去已知/未知边界）。给贪心一个公道：信息增益常具\emph{次模性}（边际递减——你已看清大半，再看一眼新信息就少了），对单调 + 次模目标贪心能达最优的 $(1-1/e)\approx0.63$（Nemhauser--Wolsey--Fisher 界），故"选哪些测量点"的集合选择上贪心有可证保证。但完整 Active SLAM 还含运动代价与序贯性（非纯次模），仍需 POMDP 多步前瞻处理"现在收益小、但解锁未来大回环"的情况。

\begin{pitfall}[概念误区：以为信息增益是唯一奖励]
只用信息增益当奖励，机器人会无限探索、永不完成任务（或不回家、不到目标）。Active SLAM 的奖励是\emph{多项加权}——信息增益 + 覆盖 $-$ 碰撞 $-$ 运动代价（$-$ 到目标距离，若有任务目标）。按任务配权重（纯探索可重信息增益；带目标的导航要加到达项）。另注：高维协方差直接求 $\det$ 会浮点下溢/溢出，应用 $\log\det\boldsymbol{\Sigma}=\sum\log\lambda_i$（对数域，数值稳定）。
\end{pitfall}

\begin{note}
\textbf{神经世界模型 $=$ 摊销 belief.}\;当观测是图像/点云这类高维连续信号时，粒子滤波撞维数灾难、符号模型也写不出来。神经世界模型（如 DreamerV3 的循环状态空间模型 RSSM）用神经网络从高维观测里学一个紧凑的隐空间 belief 表示 + 隐空间动力学——它的隐状态 $(h_t,z_t)$ 就是一个\emph{摊销的 belief}：后验 encoder $=$ belief update 的观测校正、动力学 predictor $=$ 运动预测、隐空间 actor-critic $=$ belief 空间策略优化。DESPOT 每周期\emph{现算}一棵树（贵但可解释），RSSM 把"belief update + 规划"\emph{摊销进网络权重}（快但黑箱）——两者求解同一个问题（belief 空间 MDP），合流点是 neural-guided 树搜索（学习注入先验、搜索保留前瞻与保证）。\pz{此处为跨领域指针，神经世界模型/分布式 RL 的深讲属于学习部，不在此展开长推导。}
\end{note}

\begin{practice}
\begin{enumerate}
  \item 实现 LQG-MP 的核心：给两条候选路径（一条短但靠障碍、一条长但远离），用 LQG 沿路径传播协方差、算各自碰撞概率，选更可靠的。验证"短路径不一定更优"。
  \item 在 2D 占据栅格世界实现最简 Active SLAM（栅格 + 位姿不确定作 belief、信息增益作奖励、贪心选信息增益最大方向），对比贪心主动探索 vs 随机游走的地图覆盖速度。
  \item 论证为什么信息增益奖励是 POMDP 独有、MDP 里不存在的（从"MDP 假设状态可观测"出发），并说明 QMDP 为何在 Active SLAM 上失败。
\end{enumerate}
\end{practice}

% ════════════════════════════════════════════════════════════════
\section{风险敏感规划：一致性风险测度与 CVaR}
\label{sec:risk-coherent}

\paragraph{动机：当"平均最优"不够。}
机会约束只看"违反概率有没有超过 $\delta$"，完全不管"一旦违反、有多惨"。但机器人在不确定环境里，每个方案的后果是一个\emph{随机变量}（代价分布）。要在方案之间排序、写进优化目标，必须把随机变量压成一个标量——这个映射叫\textbf{风险度量} $\rho$。\emph{你怎么压（用什么 $\rho$），就编码了你怎么看待不确定性下的好坏}。最自然的选择是期望 $\rho(Z)=\mathbb{E}[Z]$，但期望把罕见的大灾难"平均掉"了（一条路 97\% 概率 10 分钟、3\% 概率 120 分钟，期望仅 $\approx13.3$ 分钟，可那 3\% 的两小时卡死会反复发生）。

\paragraph{反面：用"不讲道理"的度量会怎样。}
第二个候选是 \textbf{VaR}（Value-at-Risk）：$\mathrm{VaR}_\alpha(Z)$ 是代价分布的 $\alpha$ 分位数（"有 $\alpha$ 的把握代价不超过这个值"）。但 VaR 有\emph{结构性}缺陷。设机器人同时执行两个独立子任务 A、B，各以 4\% 概率损失 100。各自 $\Pr(\cdot>0)=4\%<5\%$，故 $\mathrm{VaR}_{0.95}(A)=\mathrm{VaR}_{0.95}(B)=0$。但合起来：至少一个失败的概率 $1-0.96^2\approx7.8\%>5\%$，于是
\begin{equation}
  \mathrm{VaR}_{0.95}(A+B)\approx100\ >\ \mathrm{VaR}_{0.95}(A)+\mathrm{VaR}_{0.95}(B)=0.
  \label{eq:risk-var-counter}
\end{equation}
\textbf{把两个任务合起来，VaR 反而暴涨}——违反了"合并不应凭空制造风险"的直觉，会诱导机器人把任务拆开单独做（即使物理上没理由）。\pz{存疑：源给出对应 CVaR 数值 $\mathrm{CVaR}_{0.95}(A)=\mathrm{CVaR}_{0.95}(B)=0.04\times100/0.05=80$、$\mathrm{CVaR}_{0.95}(A+B)=(0.0016\times200+0.0484\times100)/0.05=103.2$，已逐步核对自洽（注意须用 RU 公式而非"对 $\ge$VaR 样本取平均"的朴素估计，后者因原子算错）。}

\subsection{Artzner 一致性风险测度的四公理}
\label{sec:risk-axioms}

\begin{definition}[一致性风险测度]\label{def:risk-coherent}
Artzner--Delbaen--Eber--Heath\cite{artzner1999coherent} 提出：一个风险度量 $\rho$（从代价随机变量到实数，代价越大越坏）若满足以下四条公理，称为\textbf{一致性的}（coherent）：
\begin{enumerate}
  \item \textbf{单调性（A1）}：若 $Z\le Z'$ 逐样本成立，则 $\rho(Z)\le\rho(Z')$（严格更优的方案不被判更危险）；
  \item \textbf{平移不变性（A2）}：对常数 $c$，$\rho(Z+c)=\rho(Z)+c$（确定代价线性反映到风险，$\rho(Z)$ 可解读为"为让 $Z$ 可接受需预留多少确定裕量"）；
  \item \textbf{正齐次性（A3）}：对 $\beta\ge0$，$\rho(\beta Z)=\beta\rho(Z)$（风险随规模等比缩放、不依赖单位）；
  \item \textbf{次可加性（A4）}：$\rho(Z+Z')\le\rho(Z)+\rho(Z')$（"合并不产生额外风险"，正面编码分散化/冗余的价值）。
\end{enumerate}
\end{definition}

A4 是四条里最深刻、也是排除 VaR 的那一条——\cref{eq:risk-var-counter} 的"合并风险暴涨"正是违反它。A3 + A4 合起来叫\emph{次可线性}，并推出\emph{凸性} $\rho(\lambda Z+(1-\lambda)Z')\le\lambda\rho(Z)+(1-\lambda)\rho(Z')$——凸性对优化至关重要（保证"把风险写进目标/约束后问题仍凸、可高效全局求解"），\cref{sec:risk-embed} 把 CVaR 嵌进 MPC 时会直接兑现。期望满足全部四条（最简单的一致性度量），但它把尾部平均掉了；VaR 违反 A4 且非凸。

\begin{note}
\textbf{机器人特有的两条公理.}\;Majumdar--Pavone\cite{majumdar2017risk} 论证机器人风险度量还应满足：\emph{A5 共单调可加性}——若 $Z,Z'$ 共单调（同一底层随机源的两个单调函数），则 $\rho(Z+Z')=\rho(Z)+\rho(Z')$（取等、不打折），堵住"虚假分散化折扣"（两个共用同一故障源的"冗余"模块不应被当真正冗余）；\emph{A6 律不变性}——同分布则同风险（风险只取决于分布，可只从样本估计）。满足 A1--A6 的恰好是\emph{畸变风险度量} $\rho_g(Z)=\int_0^1 F_Z^{-1}(\tau)\,\mathrm{d}g(\tau)$（$g$ 是单调畸变函数）——选风险度量 $=$ 选一条畸变曲线，CVaR 是 $g(\tau)=\min(\tfrac\tau\alpha,1)$ 的特例。\pz{存疑：表示定理（Kusuoka/Föllmer--Schied 一脉）的精确表述与 CVaR 对应的畸变函数约定，不同文献因"代价 vs 收益"方向略有等价变体，联网核对。}
\end{note}

\subsection{CVaR 与 Rockafellar--Uryasev 变分公式}
\label{sec:risk-cvar}

\begin{definition}[CVaR]\label{def:risk-cvar}
\textbf{CVaR}（Conditional Value-at-Risk，条件风险价值；亦称 Expected Shortfall）不看分位点，而看分位点以上整个尾部的平均：
\begin{equation}
  \mathrm{CVaR}_\alpha(Z)=\mathbb{E}\big[\,Z\mid Z\ge\mathrm{VaR}_\alpha(Z)\,\big]\quad(\text{连续分布下的直觉形式}).
  \label{eq:risk-cvar-intuitive}
\end{equation}
直觉：$\mathrm{CVaR}_{0.95}$ $=$ "最坏 5\% 情况的平均代价"。可证 CVaR 满足全部四条公理（\cref{def:risk-coherent}），是一致性风险测度。
\end{definition}

\begin{insight}
VaR 看"门槛在哪"、CVaR 看"门槛后面有多惨"——一字之差，公理性质天壤之别。把损失分布想成一条尾巴：VaR 是"尾巴从哪个坐标开始"的那个\emph{点}；CVaR 是"那条尾巴覆盖区域的平均高度"。两个方案可能 VaR 相同（尾巴起点一样），但一个尾巴又细又短（轻微超支）、另一个又粗又长（灾难性超支）——VaR 说它们一样危险，CVaR 一眼看出后者差得多。数学后果是决定性的：看"点"的 VaR 因分位点突跳而违反次可加；看"尾部平均"的 CVaR 因积分光滑而满足全部公理。
\end{insight}

\cref{eq:risk-cvar-intuitive} 虽直观，但\emph{依赖先知道 $\mathrm{VaR}_\alpha$}（先算分位点才能对超过它的部分取平均），而分位数关于决策变量\emph{非凸、非光滑}——嵌进优化会变成"内层嵌着非凸非光滑子问题"的双层结构，求解器罢工。Rockafellar--Uryasev\cite{rockafellar2000optimization} 给出了一个出人意料的解法。

\begin{theorem}[Rockafellar--Uryasev 变分公式]\label{thm:risk-ru}
CVaR 可写成对一个辅助标量 $z$ 求最小：
\begin{equation}
  \mathrm{CVaR}_\alpha(X)=\min_{z\in\mathbb{R}}\ \Big\{\,z+\frac{1}{1-\alpha}\,\mathbb{E}\big[(X-z)_+\big]\,\Big\},
  \label{eq:risk-ru}
\end{equation}
其中 $(X-z)_+=\max(0,X-z)$。右边不需要预先知道 VaR；取到最小的 $z^*$ 恰好就是 $\mathrm{VaR}_\alpha$，最小值恰好就是 $\mathrm{CVaR}_\alpha$，且目标关于 $z$ 凸。
\end{theorem}
\begin{derivation}
记 $F(z)=z+\frac{1}{1-\alpha}\mathbb{E}[(X-z)_+]$。利用 $\frac{\mathrm{d}}{\mathrm{d}z}\mathbb{E}[(X-z)_+]=-\Pr(X>z)$（$z$ 增大一点，超过 $z$ 的样本各减少同样一点，期望减少速率正是"有多大比例样本超过 $z$"），得
\[
  \frac{\mathrm{d}F}{\mathrm{d}z}=1-\frac{1}{1-\alpha}\Pr(X>z).
\]
令其为 0：$\Pr(X>z)=1-\alpha$，即 $\Pr(X\le z)=\alpha$——正是 $\alpha$ 分位数的定义，故 $z^*=\mathrm{VaR}_\alpha$。代回并整理（$\mathbb{E}[(X-z^*)_+]=(1-\alpha)\mathbb{E}[X-\mathrm{VaR}_\alpha\mid X\ge\mathrm{VaR}_\alpha]$）：
\[
  F(z^*)=\mathrm{VaR}_\alpha+\mathbb{E}[X-\mathrm{VaR}_\alpha\mid X\ge\mathrm{VaR}_\alpha]=\mathbb{E}[X\mid X\ge\mathrm{VaR}_\alpha]=\mathrm{CVaR}_\alpha.
\]
$(X-z)_+$ 对 $z$ 凸、期望保凸、加线性项仍凸，故"令导数为 0"找到的是全局最小。
\end{derivation}

\paragraph{从公式到线性规划。}
变分公式里的期望用样本平均近似（SAA），有 $N$ 个损失样本 $\{x_1,\dots,x_N\}$：
\[
  \mathrm{CVaR}_\alpha(X)\approx\min_{z}\ \Big\{z+\frac{1}{(1-\alpha)N}\sum_{i=1}^N(x_i-z)_+\Big\}.
\]
用辅助变量 $u_i\ge0$、$u_i\ge x_i-z$ 处理非光滑的 $(x_i-z)_+$（最小化时 $u_i$ 自动被压到 $\max(0,x_i-z)$），CVaR 计算变成标准\textbf{线性规划}：
\begin{equation}
  \min_{z,\,u_1,\dots,u_N}\ z+\frac{1}{(1-\alpha)N}\sum_{i=1}^N u_i
  \quad\text{s.t.}\quad u_i\ge x_i-z,\ \ u_i\ge0,\ \ i=1,\dots,N.
  \label{eq:risk-cvar-lp}
\end{equation}

\begin{insight}
要把 $\max(0,\text{表达式})$ 放进最小化问题，就引入一个辅助变量、给它"$\ge$ 表达式"和"$\ge0$"两个下界，最小化会自动把它压到 $\max$。这个"辅助变量 + 线性约束"的小技巧，是 CVaR 能进入实时机器人优化的全部秘密：若损失是决策变量 $\theta$ 的函数 $x_i=L(\theta,\omega_i)$ 且 $L$ 对 $\theta$ 凸，则约束 $u_i\ge L(\theta,\omega_i)-z$ 关于 $(\theta,z,u)$ 仍凸，整个"优化 $\theta$ 使 CVaR 最小"保持凸——这正是 \cref{sec:risk-embed} 把 CVaR 嵌进 MPC 的基础。
\end{insight}

\begin{codebox}[Python]{CVaR 的三种算法（朴素 / RU / LP）}
import numpy as np
from scipy.optimize import linprog

def cvar_naive(losses, alpha):
    # naive: mean of samples >= VaR. NOT robust for atomic distributions.
    var = np.quantile(losses, alpha)
    tail = losses[losses >= var]
    return tail.mean() if len(tail) else var

def cvar_ru(losses, alpha):
    # Rockafellar-Uryasev: z + mean((x-z)+)/(1-alpha), z = VaR. Robust.
    z = np.quantile(losses, alpha)
    return z + np.mean(np.maximum(losses - z, 0.0)) / (1.0 - alpha)

def cvar_lp(losses, alpha):
    # SAA-LP: min z + 1/((1-a)N) sum u_i  s.t. u_i >= x_i - z, u_i >= 0
    n = len(losses)
    c = np.concatenate([[1.0], np.full(n, 1.0 / ((1 - alpha) * n))])
    A_ub = np.zeros((n, 1 + n)); b_ub = np.zeros(n)
    for i in range(n):                       # u_i >= x_i - z  ->  -z - u_i <= -x_i
        A_ub[i, 0] = -1.0; A_ub[i, 1 + i] = -1.0; b_ub[i] = -losses[i]
    bounds = [(None, None)] + [(0, None)] * n
    res = linprog(c, A_ub=A_ub, b_ub=b_ub, bounds=bounds, method='highs')
    return res.fun, res.x[0]
\end{codebox}

\begin{pitfall}[编程陷阱：对有原子的损失用 mean(x[x >= VaR]) 估 CVaR]
对二值代价（碰/不碰）或离散损失，"对超过 VaR 的样本取平均"会系统性偏离真实 CVaR，甚至让你误判次可加性。例如两个独立"4\% 概率损失 100"的 $A,B$，朴素估计算出 $\mathrm{CVaR}(A)\approx4$（理论应为 $80$）。根因：分布在 VaR 处有原子（一团样本恰好等于分位点），"$\ge$ VaR 的样本"占比不等于 $1-\alpha$。正确做法：用 RU 公式 \cref{eq:risk-ru} 或解 SAA-LP \cref{eq:risk-cvar-lp}（对原子稳健）。自检：对已知解析 CVaR 的离散分布（如上面的 $A$，理论 $\mathrm{CVaR}_{0.95}=4/0.05=80$）测试，对不上就说明踩了原子坑。
\end{pitfall}

\begin{note}
\textbf{CVaR 的对偶面孔与分布鲁棒.}\;一致性风险测度有一般对偶表示 $\rho(X)=\max_{Q\in\mathcal{Q}}\mathbb{E}_Q[X]$（在一族概率测度上取最坏期望）。对 CVaR，$\mathcal{Q}$ 是所有"密度不超过 $\frac{1}{1-\alpha}$ 倍原测度"的概率测度——把全部"可放大额度"压到 $X$ 最坏的 $1-\alpha$ 尾部，恰好得到尾部平均。把这个测度族再扩大（让 $\mathcal{Q}$ 是"离名义分布不超过某距离的所有分布"），就得\emph{分布鲁棒 CVaR}（DR-CVaR）——在分布本身不确定时仍能给尾部保证，是 \cref{sec:cc-spectrum} 的 DRO 在风险度量上的兑现。\pz{控制部深讲：分布鲁棒 MPC (DR-MPC) 的 Wasserstein 球构造与求解，留待控制部鲁棒/随机 MPC 章。}
\end{note}

\begin{practice}
\begin{enumerate}
  \item 实现 RU 的 CVaR-LP \cref{eq:risk-cvar-lp}：对标准正态与重尾标准化 $t$（自由度 3）样本各算 $\mathrm{CVaR}_{0.95}$，对比 CVaR 与 VaR 的差异（重尾下 CVaR/VaR 比值显著更大），并验证 LP 解出的 $z^*$ 等于经验 VaR。
  \item 构造离散损失（4\% 概率取 100，否则 0），分别用朴素估计与 RU 公式算 $\mathrm{CVaR}_{0.95}$，验证朴素偏离理论值 80、RU 给出正确的 80。
  \item 论证变分公式 \cref{eq:risk-ru} 把"先算分位数再取尾部平均"的非凸两步变成凸的一步最小化，为什么这对"把 CVaR 写进 MPC"如此关键。
\end{enumerate}
\end{practice}

% ════════════════════════════════════════════════════════════════
\section{时间一致性、嵌套 CVaR 与嵌入安全工具}
\label{sec:risk-nested}

\paragraph{动机：单步 CVaR 好用，多步 CVaR 却会"自己打自己脸"。}
\cref{sec:risk-cvar} 在\emph{一次性决策}里用 CVaR。但机器人决策是\emph{序贯}的，我们关心整条轨迹累积代价 $L=c_0+\gamma c_1+\gamma^2c_2+\cdots$ 的尾部风险。最自然的想法是直接最小化 $\mathrm{CVaR}_\alpha(L)$——这叫\textbf{静态 CVaR}。但它有个深刻陷阱：\emph{静态 CVaR 是时间不一致的}。你在 $t=0$ 按"最小化整条轨迹 CVaR"定下的最优策略，走到 $t=1$、获得新信息后重新优化，会得到一个\emph{与原计划冲突}的策略——昨天的最优计划，今天的你想推翻。这对动态规划（\cref{thm:plan-bellman}）是致命的：DP 假设"未来的最优子策略可独立求解、再拼回整体最优"（最优性原理），时间不一致直接摧毁这个前提。

\begin{insight}
时间不一致不是 CVaR 的毛病，而是"在序贯决策里如何施加风险"的根本张力。问题出在把"对整条随机轨迹的一次性风险评估"硬塞进"逐步展开的序贯决策"这个错配上——任何"对全程总量施加一次的非线性风险泛函"都会面临时间不一致，因为"全程尾部"的定义会随信息更新而漂移。要让风险与 DP 兼容，就必须把风险的施加也\emph{序贯化}。这呼应一个更普遍的原理：当你想用动态规划求解时，目标函数的结构必须和"逐步递推"相容——期望天然相容（线性、可分解），非线性风险泛函必须改造成嵌套形式才相容。
\end{insight}

\subsection{嵌套 CVaR 与 risk-averse Bellman 方程}
\label{sec:risk-bellman}

Ruszczyński\cite{ruszczynski2010risk} 的修复思路：不在 $t=0$ 对"全程总代价"施加一次 CVaR，而是\emph{在每一步都施加一次"一步条件 CVaR"，层层嵌套}。

\begin{theorem}[嵌套 CVaR 的 risk-averse Bellman 方程]\label{thm:risk-nested}
把目标从静态的 $\mathrm{CVaR}_\alpha(\sum_t\gamma^t c_t)$ 换成嵌套形式，定义 risk-averse 值函数 $V(s)$，它满足
\begin{equation}
  V(s)=\min_{a}\ \Big\{\,c(s,a)+\gamma\,\mathrm{CVaR}_\alpha^{s}\big[\,V(s')\,\big]\,\Big\},
  \label{eq:risk-bellman}
\end{equation}
其中 $s'$ 是从 $(s,a)$ 转移到的随机下一状态，$\mathrm{CVaR}_\alpha^{s}[\cdot]$ 是对这个随机 $s'$ 取的一步条件 CVaR。它和标准 Bellman 方程几乎一样，\emph{唯一改动是把"对下一状态取期望 $\mathbb{E}_{s'}[V(s')]$"换成"取一步 CVaR"}。该构造满足：(1) 时间一致性自动成立；(2) risk-averse Bellman 方程有唯一解；(3) 一步 CVaR 算子是 $\gamma$-压缩映射（CVaR 满足单调性 + 平移不变正好保证），故值迭代/策略迭代收敛。
\end{theorem}

\paragraph{为什么修复了时间一致性。}
嵌套形式天然\emph{递归}——它的结构本身就长成 Bellman 方程的样子。直觉上，标准 Bellman 假设"下一状态按其概率平均"，risk-averse Bellman 则悲观地假设"下一状态偏向最坏 $1-\alpha$ 尾部"。因为这个改动只动了"如何聚合下一状态的随机性"、没动递归结构，求解机器（值迭代、策略迭代）几乎原样可用——这是嵌套形式相比静态形式的核心工程优势。实际算这个一步 CVaR 很直接：下一状态 $s'$ 按 $P(s'\mid s,a)$ 分布、$V(s')$ 是一组带概率的值（一个离散随机变量），用 \cref{thm:risk-ru} 的 RU 公式或小 LP \cref{eq:risk-cvar-lp} 算即可。

\cref{tab:risk-static-nested} 总结静态 vs 嵌套这个必须记住的区分。

\begin{table}[H]\centering\small
\caption{静态 CVaR vs 嵌套 CVaR}
\label{tab:risk-static-nested}
\begin{tabular}{lll}
\toprule
 & 静态 CVaR & 嵌套 CVaR（Markov）\\
\midrule
施加方式 & 对全程总代价施加\emph{一次} & 每步施加\emph{一步条件} CVaR，层层嵌套 \\
时间一致性 & 不一致（尾部定义随信息漂移）& 一致（每步只看下一步）\\
能否用 Bellman/DP & 不能直接递推 & risk-averse Bellman 成立，VI/PI 收敛 \\
保守程度 & 相对不那么保守 & 通常更保守（每步担忧 + 累积）\\
求解 & 需特殊处理（如增广 VaR 阈值进状态）& 像标准 MDP，把期望换成一步 CVaR \\
\bottomrule
\end{tabular}
\end{table}

实践含义：\textbf{要用 DP/RL 的递推机器求解多步风险敏感问题，几乎总是用嵌套 CVaR}（它和 Bellman 兼容）；若坚持"全程总代价 CVaR"的精确语义（某些安全认证场景），就得接受静态 CVaR 的不可递推性、用别的手段（如把 VaR 阈值增广进状态）。两者不可混为一谈——这也是文献常见的混淆源。值得一提：用分布式 RL（学整个回报分布，再按 CVaR 贪心）做"每步按 CVaR 选动作"得到的是\emph{动态/迭代 CVaR}（接近嵌套），\emph{不是}静态 CVaR\cite{lim2022distributional}。\pz{此处为指针，分布式 RL（C51/QR-DQN/IQN 学回报分布 $\to$ 任意畸变风险）的深讲属学习部，不在此展开。}

\begin{pitfall}[思维陷阱：以为时间不一致"只是理论洁癖"]
明知静态 CVaR 时间不一致，仍每步重新优化静态 CVaR 当没事——结果机器人行为抖动、计划反复推翻、甚至闭环里绕圈。时间不一致有\emph{真实的闭环后果}：每步重优化给出互相矛盾的计划，叠加执行就是抖动/振荡/次优。正确做法：序贯闭环里要么用嵌套 CVaR（时间一致、可放心每步重规划），要么用静态 CVaR 但\emph{承诺执行}整条规划（不中途重优化）并配增广状态正确求解。别"用静态 CVaR 的目标 + 每步重优化的执行"这种自相矛盾的组合。
\end{pitfall}

\subsection{把 CVaR 嵌进 MPC 与安全滤波}
\label{sec:risk-embed}

\paragraph{动机：确定性的安全条件遇到随机扰动就不够用。}
真实机器人状态转移是 $\mathbf{x}_{t+1}=f(\mathbf{x}_t,\mathbf{u}_t)+\mathbf{w}_t$，安全裕量函数 $h(\mathbf{x}_{t+1})$ 成了随机变量。该要求什么安全条件？\cref{tab:risk-safety} 给出几个层次，正好对应一条"安全谱"。

\begin{table}[H]\centering\small
\caption{随机扰动下安全条件的层次（安全谱）}
\label{tab:risk-safety}
\begin{tabular}{llll}
\toprule
安全条件 & 要求 & 保守程度 & 无界扰动下可行性 \\
\midrule
期望安全 & $\mathbb{E}[h(\mathbf{x}_{t+1})]\ge0$ & 最弱（对尾部失明）& 好 \\
机会约束安全 & $\Pr(h(\mathbf{x}_{t+1})\ge0)\ge1-\delta$ & 弱（对严重程度失明）& 好 \\
\textbf{CVaR 安全} & $\mathrm{CVaR}_\alpha[-h(\mathbf{x}_{t+1})]\le0$ & 中（对尾部敏感）& \textbf{可行} \\
最坏情况安全 & $\forall \mathbf{w}:\ h(\mathbf{x}_{t+1})\ge0$ & 最强 & 无界扰动下\emph{不可行} \\
\bottomrule
\end{tabular}
\end{table}

\paragraph{CVaR-MPC。}
把 \cref{thm:risk-ru} 的辅助变量嵌进 MPC：对每个预测步、对随机扰动采 $N$ 个场景，用 SAA-LP 把"碰撞代价的 $\mathrm{CVaR}_\alpha\le$ 阈值 $\kappa$"写成辅助变量 + 线性约束。记预测时域 $H$、决策为控制序列 $\mathbf{u}_{0:H-1}$、$N$ 条扰动场景 $\{\omega^{(j)}\}$：
\begin{equation}
\begin{aligned}
\min_{\mathbf{u}_{0:H-1},\,z,\,\{v_j\}}\quad & \sum_{t=0}^{H-1}\big(\|\mathbf{x}_t-\mathbf{x}_{\mathrm{ref}}\|_{\mathbf{Q}}^2+\|\mathbf{u}_t\|_{\mathbf{R}}^2\big)\\
\text{s.t.}\quad & \mathbf{x}_{t+1}=f(\mathbf{x}_t,\mathbf{u}_t),\\
& z+\frac{1}{(1-\alpha)N}\sum_{j=1}^N v_j\le\kappa,\\
& v_j\ge c_{\mathrm{collision}}\big(\mathbf{x}_{0:H}^{(j)}\big)-z,\quad v_j\ge0,\quad j=1,\dots,N,
\end{aligned}
\label{eq:risk-cvar-mpc}
\end{equation}
其中后两行就是 \cref{eq:risk-cvar-lp} 的 CVaR-LP 形式（$z$ 自动解到 VaR、$v_j$ 被压到 $(c^{(j)}-z)_+$）。CVaR 带来的额外规模\emph{只随场景数 $N$ 线性增长}（$1+N$ 个变量、$N+1$ 行约束），且\emph{不破坏凸性}——所以整体仍是凸 QP，OSQP/acados 毫秒级可解。把 $\alpha\to0$ 退化回标称 MPC，把 CVaR 约束换成 $\Pr$ 约束就是 CC-MPC——\emph{同一个 QP 骨架，换那几行约束就在安全谱上滑动}。

\begin{insight}
CVaR-CBF/MPC 不是新算法，而是给安全骨架换了一个"风险态度旋钮"。控制障碍函数（CBF）安全滤波和 MPC 的骨架（把标称控制投影到安全集 / 滚动优化）是控制部就建好的。CVaR 版的贡献不在骨架，而在\emph{安全条件的写法}——把"$h\ge0$"替换成 CVaR 形式。这个替换可行且优雅，全靠 \cref{thm:risk-ru}：它把 CVaR 写成凸的辅助变量形式，于是"CVaR 安全条件"和"确定性安全条件"一样都是 QP 里的凸约束，不破坏可解性。掌握 CVaR-CBF/MPC 的关键不是学新优化技巧，而是理解"安全条件可以是确定性/概率/尾部平均/最坏情况的不同写法，它们在同一个 QP 框架里可互换"。
\end{insight}

\begin{note}
\textbf{CVaR-CBF.}\;Ahmadi--Xiong--Ames\cite{ahmadi2021cvar} 把 CBF 的前向不变条件用 CVaR 表述：把确定性的 $h(\mathbf{x}_{t+1})\ge(1-\alpha_{\mathrm{cbf}})h(\mathbf{x}_t)$ 升级为对随机的裕量损失施加 CVaR，要求 $\mathrm{CVaR}_\alpha[h(\mathbf{x}_t)-h(\mathbf{x}_{t+1})]\le\alpha_{\mathrm{cbf}}h(\mathbf{x}_t)$（"裕量下降的最坏 $\alpha$ 尾部的平均"不超标），用 RU 辅助变量写成凸约束嵌进安全滤波 QP，在 Cassie 双足机器人上验证。\pz{控制部深讲：CBF 的前向不变性理论、类 K 函数、CBF-QP 安全滤波的完整推导，留待控制部安全控制章；此处只给"安全条件换写法"的规划/决策视角。注意两个 $\alpha$ 不同——$\alpha_{\mathrm{cbf}}$ 是 CBF 不变性速率、$\alpha$ 是 CVaR 尾部水平。}
\end{note}

\begin{pitfall}[思维陷阱：以为 CVaR 安全条件能提供和鲁棒方法一样的"硬保证"]
用 CVaR-CBF 后就当机器人"绝对安全"了——错。CVaR 控制的是\emph{尾部的平均}，不是\emph{最坏情况的绝对界}。$\mathrm{CVaR}_\alpha[-h]\le0$ 保证"最坏 $\alpha$ 尾部平均而言安全"，但尾部内部仍可能有个别更极端的违反（均值达标不代表无一例外）。它给的是风险意义的保证，不是"对所有扰动都安全"的硬保证。正确做法：清楚 CVaR 保证的语义边界——它在"鲁棒太保守/不可行"时提供可行且对尾部敏感的中间方案，但若任务要求绝对硬安全（且扰动有界），仍需鲁棒方法。另：CVaR-MPC 采样应采\emph{整条轨迹}的扰动序列（保留时间相关性），别对每步独立采样。
\end{pitfall}

\begin{practice}
\begin{enumerate}
  \item 构造一个最简两步 MDP，用静态 $\mathrm{CVaR}$ 算 $t=0$ 的最优完整计划，再模拟"走到 $t=1$、观察事件后重新优化"，验证它与原计划冲突（时间不一致）；改用嵌套 CVaR，验证 $t=0$ 计划与 $t=1$ 重优化一致。
  \item 在小格子世界 MDP（带随机转移）上实现 risk-averse 值迭代（把 $\mathbb{E}[V(s')]$ 换成一步 $\mathrm{CVaR}_\alpha[V(s')]$），对比 $\alpha\to0$（退化回标准 VI）与 $\alpha=0.9$（风险厌恶）的策略差异，验证 VI 收敛。
  \item 对 2D 点机器人 + 混合高斯扰动（重尾），实现 CVaR-MPC \cref{eq:risk-cvar-mpc} 与机会约束 MPC，蒙特卡洛对比，验证 CVaR-MPC 更好地控制了尾部碰撞代价（机会约束达标但偶发重创）。
\end{enumerate}
\end{practice}

% ════════════════════════════════════════════════════════════════
\section{常见误解汇总}
\label{sec:plan-chance-misc}

\begin{enumerate}
  \item \textbf{"每步满足 $\delta$ 就等于整条轨迹满足 $\delta$"}——错。ICC $\ne$ JCC，差一个 union bound（\cref{thm:cc-boole}）；$N$ 步每步 $\delta$，全程约 $N\delta$。要 JCC $=\delta$ 须 $\sum\delta_{j,k}\le\delta$。
  \item \textbf{"IRA 能消除 Boole 的保守度"}——错。IRA 优化的是"在 Boole 框架内怎么分"，消不掉 union bound 自身的保守度（\cref{sec:cc-ira-struct}）。
  \item \textbf{"机会约束就是鲁棒约束的概率版，$\delta\to0$ 即鲁棒"}——部分对。$\delta\to0$ 时高斯收紧 $\to\infty$（比鲁棒还保守）；只有把高斯截断成有界集，$\delta\to0$ 才退回 tube（\cref{sec:cc-vs-tube}）。
  \item \textbf{"belief 就是滤波器估出的那个状态，取均值用就行"}——错。belief 是整个分布，塌缩成点丢掉"我有多不确定"这个 POMDP 决策的核心依据（\cref{sec:pomdp-belief}）。
  \item \textbf{"$\boldsymbol{\alpha}$-vector 和 belief 都是 $|S|$ 维向量，是一回事"}——错。belief 是概率分布（非负和 1），$\boldsymbol{\alpha}$-vector 是价值向量（可负、不归一），通过内积配对（\cref{sec:pomdp-curse}）。
  \item \textbf{"$\boldsymbol{\alpha}$-vector 爆炸是实现不好，换算法就行"}——错。这是问题固有复杂度（两个 curse，PSPACE-hard），剪枝只延缓不消除（\cref{sec:pomdp-curse}）。
  \item \textbf{"点基方法给出 $V^*$ 的上界"}——错，给的是\emph{下界}（每个 $\boldsymbol{\alpha}$ 对应可执行计划，价值 $\le$ 最优；\cref{sec:pomdp-pbvi}）。
  \item \textbf{"DESPOT 的场景是采样的状态轨迹"}——错。场景 $=$ 起始状态 + 随机数串；同场景 + 不同动作 $=$ 不同轨迹（\cref{def:pomdp-despot}）。
  \item \textbf{"VaR 只是没 CVaR 好，调高 $\alpha$ 就能补救"}——错。VaR 的缺陷是结构性的（违反次可加），无论 $\alpha$ 多高都只是一个分位点、对尾部内部失明（\cref{sec:risk-axioms}）。
  \item \textbf{"对有原子的损失用 mean(x[x>=VaR]) 估 CVaR"}——错。原子处会算错，须用 RU 公式或 LP（\cref{thm:risk-ru}）。
  \item \textbf{"静态 CVaR 和嵌套 CVaR 在多步问题里是同一个东西"}——错。静态对全程施加一次、时间不一致；嵌套每步施加、时间一致（\cref{tab:risk-static-nested}）。
  \item \textbf{"CVaR 安全条件提供和鲁棒一样的硬保证"}——错。CVaR 控制尾部\emph{平均}、非最坏情况绝对界（\cref{sec:risk-embed}）。
\end{enumerate}

% ════════════════════════════════════════════════════════════════
\section*{本章小结}
\label{sec:plan-chance-summary}

本章把"不确定性下的规划与决策"沿三支柱讲清：扰动是概率的（机会约束）、状态看不清的（POMDP）、要管尾部严重程度的（风险敏感）。

\subsection*{符号表}
\begin{table}[H]\centering\small
\caption{本章主要符号}
\label{tab:plan-chance-symbols}
\begin{tabular}{lll}
\toprule
符号 & 含义 & 首次出现 \\
\midrule
$\delta$ & 风险预算（允许违反概率）、置信 $1-\delta$ & \cref{eq:cc-def} \\
$\Phi,\Phi^{-1}$ & 标准正态 CDF 及其逆（分位函数）& \cref{thm:cc-gauss} \\
$\bar{\mathbf{x}},\boldsymbol{\Sigma},\mathbf{P}_k$ & 状态均值、协方差、沿轨迹协方差 & \cref{eq:cc-gauss-tighten,eq:cc-cov-prop} \\
$\delta_{j,k}$ & 第 $k$ 步第 $j$ 面分配的风险 & \cref{def:cc-icc} \\
$\lambda$ & 拉格朗日乘子 / 正则化系数 & \cref{eq:cc-lagrangian,eq:pomdp-rdespot} \\
$b,b(s)$ & belief（状态后验）及其分量 & \cref{eq:pomdp-belief-update} \\
$\Omega,O$ & 观测空间、观测模型 & \cref{def:pomdp-seven} \\
$\gamma$ & 折扣因子 & \cref{def:pomdp-seven} \\
$\boldsymbol{\alpha},\Gamma$ & $\alpha$-vector（价值向量）及其集合 & \cref{eq:pomdp-alpha} \\
$V^*(b)$ & belief 空间最优值函数 & \cref{eq:pomdp-bellman} \\
$\phi,K$ & 确定化场景、场景数 & \cref{def:pomdp-despot} \\
$\rho,\mathrm{CVaR}_\alpha$ & 风险度量、条件风险价值 & \cref{def:risk-coherent,def:risk-cvar} \\
$\alpha$（CVaR）& 尾部水平（看最坏 $1-\alpha$）& \cref{def:risk-cvar} \\
$z,u_i,v_j$ & RU 辅助变量（解到 VaR / 压到 $(\cdot)_+$）& \cref{eq:risk-cvar-lp,eq:risk-cvar-mpc} \\
\bottomrule
\end{tabular}
\end{table}

\subsection*{定理速查表}
\begin{table}[H]\centering\small
\caption{本章主要定理与公式速查}
\label{tab:plan-chance-thm}
\begin{tabular}{lll}
\toprule
名称 & 内容 & 出处 \\
\midrule
Boole 拆解 & $\sum\delta_{j,k}\le\delta\Rightarrow$ JCC & \cref{thm:cc-boole} \\
解析高斯收紧 & $a^\top\bar{\mathbf{x}}\le b-\Phi^{-1}(1-\delta)\sqrt{a^\top\boldsymbol{\Sigma}a}$ & \cref{thm:cc-gauss} \\
闭环协方差传播 & $\mathbf{P}_{k+1}=(\mathbf{A}-\mathbf{B}\mathbf{K})\mathbf{P}_k(\mathbf{A}-\mathbf{B}\mathbf{K})^\top+\boldsymbol{\Sigma}_w$ & \cref{eq:cc-cov-prop} \\
CS 不可解耦 & 有机会约束时均值与协方差耦合 & \cref{thm:cc-cs-coupling} \\
CC $\leftrightarrow$ Safe RL & 风险预算 $\delta$ 与乘子 $\lambda$ 对偶 & \cref{sec:cc-saferl} \\
belief 充分统计量 & POMDP $=$ belief 空间 MDP & \cref{thm:pomdp-suff} \\
Sondik PWLC & $V^*(b)=\max_{\boldsymbol{\alpha}}\boldsymbol{\alpha}\cdot b$ & \cref{thm:pomdp-pwlc} \\
DESPOT output-sensitive & regret 依最优策略大小、非状态空间 & \cref{thm:pomdp-output} \\
Rockafellar--Uryasev & $\mathrm{CVaR}_\alpha=\min_z\{z+\frac{1}{1-\alpha}\mathbb{E}[(X-z)_+]\}$ & \cref{thm:risk-ru} \\
risk-averse Bellman & $\mathbb{E}_{s'}\to\mathrm{CVaR}_\alpha^s$，时间一致、压缩 & \cref{thm:risk-nested} \\
\bottomrule
\end{tabular}
\end{table}

\subsection*{知识点总表}
\begin{table}[H]\centering\small
\caption{本章知识点与难度}
\label{tab:plan-chance-points}
\begin{tabular}{lll}
\toprule
知识点 & 关键内容 & 难度 \\
\midrule
机会约束形式化 & ICC/JCC、Boole、风险分配 & $\bigstar\bigstar$ \\
解析高斯收紧 & $\Phi^{-1}$ margin、SOCP、CC-MPC & $\bigstar\bigstar$ \\
迭代风险分配 IRA & 两层结构、与 Safe RL 对偶 & $\bigstar\bigstar\bigstar$ \\
CC-RRT / 协方差控制 & 协方差传播 + 概率剔除、CS 不可解耦 & $\bigstar\bigstar\bigstar$ \\
POMDP 与 belief & 七元组、贝叶斯滤波、充分统计量 & $\bigstar\bigstar$ \\
$\alpha$-vector / 两类灾难 & Sondik PWLC、维数/历史灾难、QMDP & $\bigstar\bigstar\bigstar$ \\
点基与在线树搜索 & SARSOP、POMCP、DESPOT/R-DESPOT & $\bigstar\bigstar\bigstar\bigstar$ \\
信念空间规划 / Active SLAM & LQG-MP/FIRM、信息增益作奖励 & $\bigstar\bigstar\bigstar$ \\
一致性风险测度 & Artzner 公理、VaR 反例 & $\bigstar\bigstar\bigstar$ \\
CVaR 与 RU 公式 & 变分公式、LP 化、嵌入 MPC & $\bigstar\bigstar\bigstar$ \\
时间一致性与嵌套 CVaR & 静态不一致、risk-averse Bellman & $\bigstar\bigstar\bigstar\bigstar$ \\
\bottomrule
\end{tabular}
\end{table}

% ════════════════════════════════════════════════════════════════
\section*{延伸阅读}
\label{sec:plan-chance-reading}

\begin{itemize}
  \item \textbf{机会约束}：Blackmore--Ono--Williams\cite{blackmore2011chance}（高斯避障 CC 的权威）；Ono--Williams\cite{ono2008iterative}（IRA 原始）；Campi--Garatti\cite{campi2008exact}（场景法样本复杂度）；Calafiore--El Ghaoui\cite{calafiore2006distributionally} 与 Nemirovski--Shapiro\cite{nemirovski2006convex}（CC 复杂度与凸近似/DRO）。
  \item \textbf{采样型与协方差控制}：Luders--Kothari--How\cite{luders2010chance}（CC-RRT）；Okamoto--Goldshtein--Tsiotras\cite{okamoto2018optimal}（协方差控制）。
  \item \textbf{POMDP}：Kaelbling--Littman--Cassandra\cite{kaelbling1998planning}（POMDP 标准化）；Sondik\cite{sondik1971optimal}（PWLC 奠基）；点基 PBVI/HSVI/SARSOP\cite{pineau2003pbvi,smith2004hsvi,kurniawati2008sarsop}；在线 POMCP/DESPOT\cite{silver2010pomcp,somani2013despot}。
  \item \textbf{信念空间与 Active SLAM}：LQG-MP\cite{vandenberg2011lqgmp}、RRBT\cite{bry2011rrbt}、FIRM\cite{aghamohammadi2014firm}、BSP-iSAM\cite{indelman2015bsp}；Active SLAM 综述\cite{placed2023survey}。
  \item \textbf{风险敏感}：Artzner et al.\cite{artzner1999coherent}（一致性公理）；Rockafellar--Uryasev\cite{rockafellar2000optimization}（CVaR 优化）；Ruszczyński\cite{ruszczynski2010risk}（嵌套 CVaR）；Majumdar--Pavone\cite{majumdar2017risk}（机器人风险公理）；Ahmadi--Xiong--Ames\cite{ahmadi2021cvar}（CVaR-CBF）。
\end{itemize}

% ════════════════════════════════════════════════════════════════
\section*{与后续章节的关系}
\label{sec:plan-chance-next}

\begin{itemize}
  \item \textbf{控制部（鲁棒/随机 MPC、安全控制）}：本章的机会约束 MPC、协方差控制、CVaR-MPC/CBF 给出了"规划/决策用途"的简述与关键式；tube MPC 的鲁棒正不变集构造、CBF 的前向不变性理论、分布鲁棒 MPC 的 Wasserstein 球、HJ 可达性等\emph{最深的推导}留待控制部专章（本章多处已用 \pz{控制部深讲} 标注）。
  \item \textbf{学习部（强化学习/世界模型）}：本章点出的 CC $\leftrightarrow$ Safe RL 的 primal--dual 对偶（\cref{sec:cc-saferl}）、分布式 RL 与任意畸变风险（\cref{sec:risk-nested}）、神经世界模型 $=$ 摊销 belief（\cref{sec:pomdp-active}）都在学习部展开。
  \item \textbf{回看 \cref{ch:planning}}：CC-RRT 是 RRT/RRT$^*$（\cref{sec:plan-rrt,sec:plan-rrtstar}）的不确定性版；POMDP 把动态规划（\cref{thm:plan-bellman}）搬到 belief 空间。
\end{itemize}

% ════════════════════════════════════════════════════════════════
\section*{故障排查}
\label{sec:plan-chance-trouble}

\begin{table}[H]\centering\small
\caption{常见故障与排查}
\label{tab:plan-chance-trouble}
\begin{tabular}{p{0.30\textwidth}p{0.30\textwidth}p{0.32\textwidth}}
\toprule
现象 & 可能原因 & 处理 \\
\midrule
CC-MPC 蒙特卡洛碰撞率远超 $\delta$ & 用了 ICC 当 JCC（未 Boole 拆）；或开环协方差发散 & 按 \cref{thm:cc-boole} 拆 JCC、$\sum\delta_{j,k}\le\delta$；用闭环 \cref{eq:cc-cov-prop} \\
CC-MPC 过保守、窄通道过不去 & 均分浪费预算；$\delta$ 取太小 & 上 IRA（\cref{sec:cc-ira}）按需分配；适当放大 $\delta$ \\
IRA 几轮后违反率超标 & 未重新归一化、$\sum\delta_{j,k}$ 漂移 & 每轮强制 \texttt{deltas *= delta/deltas.sum()} \\
CC-RRT 树长不出去 & 开环协方差发散、margin 巨大 & 用闭环 $\mathbf{A}-\mathbf{B}\mathbf{K}$；rewire 后重检 CC \\
belief 多步后和不为 1 / 全 0 & 忘归一化或归一化在校正前；下溢 & 算完 $\tilde b$ 再统一归一；大状态用对数域 \\
点估计决策在双峰 belief 上荒谬 & 把 belief 塌缩成均值 & 在整个 belief 上决策（$\boldsymbol{\alpha}$ 内积/粒子仿真）\\
精确 POMDP 求解超时 & 两个 curse（$\boldsymbol{\alpha}$ 指数爆炸）& 中等规模上 SARSOP、大规模上 DESPOT \\
DESPOT 实跑差、决策抖动 & $K$ 太小过拟合场景 & 增大 $K$ + R-DESPOT 正则化（\cref{eq:pomdp-rdespot}）\\
高维 $\det\boldsymbol{\Sigma}$ 下溢/溢出 & 直接求行列式 & 用 $\log\det=\sum\log\lambda_i$ \\
离散损失 CVaR 估计偏离理论 & 朴素 mean(x[x>=VaR]) 踩原子坑 & 用 RU 公式 \cref{eq:risk-ru} 或 LP \cref{eq:risk-cvar-lp} \\
多步 CVaR 规划行为反复横跳 & 静态 CVaR 时间不一致 & 改嵌套 CVaR（\cref{thm:risk-nested}）或承诺执行 \\
\bottomrule
\end{tabular}
\end{table}
